\subsection{算子的积分}

设 G 与 H 为 R 上的两个 Hausdorff 局部凸空间，且现在假定 F 是由 G 到 H 中的连续线性映射所成的空间 \(\mathcal{L}(\mathrm{G};\mathrm{H})\) ，赋以逐点收敛拓扑。F 上的每个连续线性形式都可以延拓为乘积空间 \(H^{G}\) 上的连续线性形式（TVS, II, §4, No. 1, 命题 2），因而可以写成 \(u \mapsto \sum_{i=1}^{n} \langle u(\mathbf{a}_{i}), \mathbf{b}_{i}^{\prime} \rangle\) ，其中诸 \(a_{i}\)（相应地，诸 \(b_{i}^{\prime}\) ）是 G 的元素（相应地，H 的对偶 \(H^{\prime}\) 的元素）。说 T 到 F 中的映射 U 是标量本质 \(\mu\)-可积的，意思是对每个 \(a \in G\) 与每个 \(b \in H^{\prime}\) ，数值函数 \(t \mapsto \langle U(t) \cdot a, b^{\prime} \rangle\) 都是本质 \(\mu\)-可积的。

命题 9. --- 设 U 是标量本质 \(\mu\)-可积的、T 到 \(\mathbf{F} = \mathcal{L}_{s}(\mathbf{G}; \mathbf{H})\) 中的映射。为使 \(\int U d\mu \in F\) ，必须且只需下列两个条件成立：

\begin{enumerate}
\def\labelenumi{\alph{enumi})}
\item
  对每个 \(\mathbf{x} \in \mathbf{G}\) ， \(\int (U(t) \cdot \mathbf{x}) d\mu(t) \in \mathrm{H}\) 。
\item
  设 \(B'\) 为 \(H'\) 的任意等度连续子集；由线性形式 \(u_{y'} : x \mapsto \int \langle U(t) \cdot x, y' \rangle d\mu(t)\) （其中 \(y'\) 跑遍 \(B'\) ）所成的集是等度连续的。
\end{enumerate}

条件 a) 与 b) 是必要的。因为，对每个 \(\mathbf{x} \in G\) ，映射 \(\widetilde{\mathbf{x}}: V \mapsto V \cdot \mathbf{x}\) （由 \(\mathcal{L}_s(\mathrm{G};\mathrm{H})\) 到 H 中）是线性的，所以可以看出（No. 1，

命题 1）：\(\widetilde{\mathbf{x}}\circ U:t\mapsto U(t)\cdot \mathbf{x}\) 是标量本质 \(\mu\)-可积的，并且

\[
S \cdot \mathbf {x} = \int (U (t) \cdot \mathbf {x}) d \mu (t),\tag{1}
\]

其中 \(S = \int U d\mu \in \mathcal{L}_s(\mathrm{G};\mathrm{H})\) 。这就证明了 a) 。此外， (1) 式也可以写成

\[
\langle S \cdot \mathbf {x}, \mathbf {y} ^ {\prime} \rangle = \int \left\langle U (t) \cdot \mathbf {x}, \mathbf {y} ^ {\prime} \right\rangle d \mu (t) = \left\langle \mathbf {x}, u _ {\mathbf {y} ^ {\prime}} \right\rangle ,\tag{2}
\]

换言之 \(^{t}S \cdot y' = u_{y'}\) 。由于 S 连续， \(^{t}S\) 把 \(H'\) 的每个等度连续子集变换成 \(G'\) 的等度连续子集，由此得 b) 。

反之，设 a) 与 b) 成立。根据 a) ，公式 (1) 定义了一个由 G 到 H 中的线性映射 S ，且对每个 \(y' \in H'\) ，该映射满足 (2) （No. 1, 命题 1）；而这时条件 b) 表明 S 连续（TVS, II, §6, No. 4, 命题 5 与命题 6, 以及 III, §3, No. 5, 命题 7），因此 \(S \in \mathcal{L}_{s}(G;H)\) 。最后，公式 (2) 证明 \(S = \int U d\mu\) 。

推论. --- 命题 9 的条件 b) 在下列两种情形的每一种中都成立：

\(1^{\circ}\) 测度 \(\mu\) 有界，且若 S 为其支集，则 \(U(S)\) 是 \(\mathcal{L}(G;H)\) 的等度连续子集。

\(2^{\circ}\) 命题 9 的条件 a) 成立，空间 G 是桶式的，且对 T 的每个紧子集 K ， \(U(\mathrm{K})\) 是 \(\mathcal{L}_s(\mathrm{G};\mathrm{H})\) 的有界子集。

先考察情形 \(1^{\circ}\) 。我们可以限于 S = T 的情形（Ch. V, §7, No. 1, 定理 1）。于是，设 \(B'\) 为 \(H'\) 的任意等度连续子集，存在一个等度连续、凸、平衡且弱闭的子集 \(A' \subset G'\) ，使得 \({}^{t}U(t) \cdot y' \in A'\) 对一切 \(y' \in B'\) 与一切 \(t \in T\) 成立（TVS, II, §6, No. 4, 命题 6）。由于 U 是标量 \(\mu\)-可积的，映射 \(t \mapsto {}^{t}U(t) \cdot y'\) （取值于 G 的对偶 \(G'\) ，后者赋有 \(\sigma(G', G)\) ）是标量 \(\mu\)-可积的，于是可以写出

\[
u _ {\mathbf {y} ^ {\prime}} = \int \left(^ {t} U (t) \cdot \mathbf {y} ^ {\prime}\right) d \mu (t).
\]

由于 \(A'\) 对 \(\sigma(G', G)\) 是凸的且紧的，第 2 目命题 5 的推论表明 \(u_{\mathbf{y}'} \in \mu(T)A'\) （对每个 \(\mathbf{y}' \in B'\) ），这就证明了我们的断言。

现在考察情形 \(2^{\circ}\) 。对每个 \(y' \in H'\) 与 T 的每个紧子集 K ，令

\[
u _ {\mathrm{K}, \mathbf {y} ^ {\prime}} = \int \varphi_ {\mathrm{K}} (t) \left(^ {t} U (t) \cdot \mathbf {y} ^ {\prime}\right) d \mu (t),
\]

它是代数对偶 \(\mathbf{G}^*\) （ \(\mathbf{G}\) 的代数对偶）中的元素。由于 \(\mathbf{G}\) 是桶式的， \(\mathcal{L}_s(\mathbf{G};\mathbf{H})\) 的每个有界子集都等度连续（TVS, III, §4, No. 2, 定理 1）；把论证的第一部分应用于函数 \(\varphi_{\mathbf{K}}U\) 和有界测度 \(\varphi_{\mathbf{K}}\cdot \mu\) ，便得 \(u_{\mathbf{K},\mathbf{y}'}\in \mathbf{G}'\) 。此外，对拓扑 \(\sigma (\mathbf{G}^{*},\mathbf{G})\) ，有 \(u_{\mathbf{y}'} = \lim_{\mathbf{K}} = u_{\mathbf{K},\mathbf{y}'}\) ，极限是沿 \(\mathbf{T}\) 的紧子集所成的递增有向集取的（Ch. V, §1, No. 3, 命题 10）。为验证命题 9 的条件 b) ，根据命题 9，只需证明线性映射 \(S\) （由 (1) 定义，把 \(\mathbf{G}\) 映到 \(\mathbf{H}\) 中）是连续的；再者， \(\mathbf{G}\) 是桶式的，故只需证明 \(S\) 在 \(\mathbf{G}\) 与 \(\mathbf{H}\) 赋以它们的弱化拓扑时是连续的（TVS, IV, §1, No. 3, 命题 7）；最后，根据 (2) ，归结为证明： \(u_{\mathbf{y}'}\in \mathbf{G}'\) 对每个 \(\mathbf{y}'\in \mathbf{H}'\) 成立 。由于 \(u_{\mathbf{y}'}\) ——对拓扑 \(\sigma (\mathbf{G}^{*},\mathbf{G})\) 而言——属于集 \(\mathbf{M}'\) （由诸 \(u_{\mathbf{K},\mathbf{y}'}\) 组成，其中 \(\mathbf{K}\) 跑遍 \(\mathbf{T}\) 的紧子集所成之集）的闭包，只需证明 \(\mathbf{M}'\) 等度连续（TVS, III, §3, No. 4, 命题 4）；又由于 \(\mathbf{G}\) 是桶式的，这等价于说：对每个 \(\mathbf{x}\in \mathbf{G}\) ，诸 \(\langle \mathbf{x},u_{\mathbf{K},\mathbf{y}'}\rangle\) 所成之集有界（TVS, III, §4, No. 2, 定理 1）。而这由关系式

\[
| \langle \mathbf {x}, u _ {\mathrm{K}, \mathbf {y} ^ {\prime}} \rangle | = \left| \int \varphi_ {\mathrm{K}} (t) \langle U (t) \cdot \mathbf {x}, \mathbf {y} ^ {\prime} \rangle d \mu (t) \right| \leqslant \int | \langle U (t) \cdot \mathbf {x}, \mathbf {y} ^ {\prime} \rangle | d \mu (t).
\]

命题 10. --- 设 U 是 T 到 F = \(\mathcal{L}_s(\mathrm{G};\mathrm{H})\) 中的映射。在下列三种情形的每一种中， U 都是标量本质 \(\mu\)-可积的，且 \(\int U d\mu \in \mathcal{L}_s(\mathrm{G};\mathrm{H})\) ：

\begin{enumerate}
\def\labelenumi{\alph{enumi})}
\item
  H 拟完备， \(\mu\) 有界，且若 S 为其支集，则 U 是 \(\mu\)-可测的且 \(U(S)\) 等度连续。
\item
  H 半自反， \(\mu\) 有界，且若 S 为其支集，则 U 是标量 \(\mu\)-可测的且 \(U(S)\) 等度连续。
\item
  H 半自反， G 桶式， U 标量本质 \(\mu\)-可积，且对 T 的每个紧子集 K ， \(U(\mathrm{K})\) 有界。
\end{enumerate}

在三种情形中， U 标量本质可积这一事实都是明显的；根据命题 9 及其推论，在每种情形只需验证命题 9 的条件 a) 。而在第一种情形，这一条件由第 2 目命题 8 得出，在后两种情形则由第 2 目命题 7 的推论得出。

\subsection{性质（GDF）}

在本目中，我们将考察具有以下性质（称为“可数闭图”性质）的局部凸空间 F ： \(^{6}\)

(GDF) 若 u 是 F 到 Banach 空间 B 中的线性映射，使得在乘积空间 \(F \times B\) 中， u 的图像 \(\Gamma\) 的点的每个收敛序列的每个极限仍属于 \(\Gamma\) ，则 u 连续。

每个 Fréchet 空间都具有性质（GDF）（TVS, I, §3, No. 3, 定理 1 的推论 5）。在附录中，我们将看到具有性质（GDF）的空间的其他例子。

命题 11. --- 每个具有性质（GDF）的 Hausdorff 局部凸空间 F 都是桶式的。

令 V 为 F 中的一个桶， q 是它的度规，它是 F 上的一个半范数；令 H 为与赋以此单个半范数所定义的拓扑的空间 F 相关联的 Hausdorff 空间。H 的完备化 \(\hat{H}\) 是一个 Banach 空间；令 \(\pi\) 为 F 到 \(\hat{H}\) 中的典范映射；我们将证明 \(\pi\) 是连续的（对 F 上的原拓扑）；这样命题就得到证明，因为 V —— \(\pi\) 下 \(\hat{H}\) 的单位球的逆像——届时将是 F 中 0 的一个邻域。为建立 \(\pi\) 的连续性，根据（GDF），只需证明 \(\pi\) 的图像在 \(F \times \hat{H}\) 中是闭的；换言之，我们必须看到：若 \(\mathfrak{F}\) 是 F 上的一个滤子，收敛于 \(x \in F\) ，且它的像 \(\pi(\mathfrak{F})\) 收敛于 \(y \in \hat{H}\) ，则 \(y = \pi(x)\) 。而今，元素 \(x'\) （跑遍极集 \(V^{\circ}\) ，即 V 在 \(F'\) 中的极集）都可以以唯一的方式延拓为 \(\hat{H}\) 上的连续线性形式（仍记作 \(x'\) ），且这些形式所成之集是 \(\hat{H}\) 的对偶的单位球；因此只需证明 \(\langle y, x' \rangle = \langle \pi(x), x' \rangle\) （对每个 \(x' \in V^{\circ}\) ）。而这由关系式

\[
\langle \mathbf {y}, \mathbf {x} ^ {\prime} \rangle = \lim _ {\mathfrak {F}} \left\langle \pi (\mathbf {z}), \mathbf {x} ^ {\prime} \right\rangle = \lim _ {\mathfrak {F}} \left\langle \mathbf {z}, \mathbf {x} ^ {\prime} \right\rangle = \left\langle \mathbf {x}, \mathbf {x} ^ {\prime} \right\rangle = \left\langle \pi (\mathbf {x}), \mathbf {x} ^ {\prime} \right\rangle .
\]

定理 1 （Gelfand--Dunford）. --- 设 F 是具有性质（GDF）的 Hausdorff 局部凸空间， \(F_{s}^{\prime}\) 是它的弱对偶。对 T 到 \(F_{s}^{\prime}\) 中的每个标量本质 \(\mu\)-可积映射 f ，积分 \(\int f d\mu\) 属于 \(F^{\prime}\) 。

回忆一下， \(F_{s}^{\prime}\) 的对偶是 F（TVS, II, §6, No. 2, 命题 3）。于是对每个 \(z \in F\) ，数值函数 \(\langle z, f \rangle\) 是本质 \(\mu\)-可积的；令 \(\theta(z)\) 为它在 \(L^{1}(\mu)\) 中的类。为证明 \(\int f d\mu \in F^{\prime}\) ，必须证明线性形式 \(z \mapsto \langle z, \int f d\mu \rangle\) 在 F 上连续；实际上，我们将证明以下更强的结果：

引理 2. --- 设 f 是 T 到 \(F_{s}^{\prime}\) 中的映射，使得对每个 \(z \in F\) ，数值函数 \(\langle z, f \rangle\) 属于 \(\overline{\mathcal{L}^{p}}(\mu) (1 \leqslant p \leqslant +\infty)\) ；令 \(\theta(z)\) 为该函数在 \(L^{p}(\mu)\) 中的类。则 \(z \mapsto \theta(z)\) 是 F 到 \(L^{p}(\mu)\) 中的连续线性映射。

根据性质（GDF），只需证明：对 F 中元素的每个收敛于 z 的序列 \((\mathbf{z}_{n})\) ，若 \((\theta(\mathbf{z}_{n}))\) 收敛于 \(u \in \mathrm{L}^{p}(\mu)\) ，则有 \(u = \theta(\mathbf{z})\) 。而今，若有必要，把序列 \((\mathbf{z}_{n})\)

换成它的一个子列，就可以假设函数序列 \(\langle \mathbf{z}_n,\mathbf{f}\rangle\) 局部几乎处处收敛到一个函数 \(h\in \overline{\mathcal{L}^p} (\mu)\) ， \(u\) 是它在 \(\mathrm{L}^p (\mu)\) 中的类（Ch. IV, §3, No. 4, 定理 3 以及 Ch. V, §1, No. 3）。由于按假设，对每个 \(t\in \mathrm{T}\) ，序列 \((\langle \mathbf{z}_n,\mathbf{f}(t)\rangle)\) 收敛于 \(\langle \mathbf{z},\mathbf{f}(t)\rangle\) ，故局部几乎处处有 \(h(t) = \langle \mathbf{z},\mathbf{f}(t)\rangle\) ，从而 \(u = \theta (\mathbf{z})\) 。

推论 1. --- 设 \(G_{i}\) ( \(1 \leqslant i \leqslant n\) ) 是 n 个具有性质（GDF）的 Hausdorff 局部凸空间， F 是 \(\prod_{i=1}^{n} G_{i}\) 上分别连续的多线性形式所成的空间，赋以逐点收敛拓扑。对 T 到 F 中的每个标量本质 \(\mu\)-可积映射 f ，有 \(\int f d\mu \in F\) 。

空间 F 与张量积 \(\bigotimes_{i=1}^{n}G_{i}\) 处于对偶，而 F 上的逐点收敛拓扑不是别的，正是拓扑 \(\sigma(F,\bigotimes_{i=1}^{n}G_{i})\) 。因此代数对偶 \(F^{\prime*}\) 是 \(\prod_{i=1}^{n}G_{i}\) 上全体多线性形式所成的空间。设 \(\mathbf{z}=(\mathbf{z}_{1},\ldots,\mathbf{z}_{n})\) 是 \(\prod_{i=1}^{n}G_{i}\) 的一个元素；对每个多线性形式 \(u\in F^{\prime*}\) ，映射 \(\mathbf{x}\mapsto\mathbf{u}(\mathbf{z}_{1},\ldots,\mathbf{z}_{i-1},\mathbf{x},\mathbf{z}_{i+1},\ldots,\mathbf{z}_{n})\) 是 \(G_{i}\) 上的一个线性形式，我们把它记作 \(\lambda_{i}(\mathbf{z})(\mathbf{u})\) ；于是我们得到线性映射 \(\lambda_{i}(\mathbf{z})\) ，它把 \(F^{\prime*}\) 映入代数对偶 \(G_{i}^{*}\) （ \(G_{i}\) 的代数对偶）中，且对拓扑 \(\sigma(F^{\prime*},\bigotimes_{i=1}^{n}G_{i})\) 与 \(\sigma(G_{i}^{*},G_{i})\) 连续。说 \(u\in F\) 就意味着：对每个指标 i 与每个 \(z\in\Pi_{i=1}^{n}G_{i}\) ，有 \(\lambda_{i}(\mathbf{z})(\mathbf{u})\in G_{i}^{\prime}\) 。于是，根据第 1 目命题 1，映射 \(\lambda_{i}(\mathbf{z})\circ\mathbf{f}\) 是标量本质 \(\mu\)-可积的、T 到 \(G_{i}^{\prime}\) （赋有拓扑 \(\sigma(G_{i}^{\prime},G_{i})\) ）中的映射，并且

\[
\int \left(\lambda_ {i} (\mathbf {z}) \circ \mathbf {f}\right) d \mu = \lambda_ {i} (\mathbf {z}) \left(\int \mathbf {f} d \mu\right).
\]

根据定理 1， \(\int (\lambda_i(\mathbf{z})\circ \mathbf{f})d\mu \in \mathrm{G}_i^\prime\) （对 \(1\leqslant i\leqslant n\) ），因此 \(\int \mathbf{f}d\mu \in \mathbf{F}\)

推论 2. --- 设 G 是具有性质（GDF）的 Hausdorff 局部凸空间， H 是半自反空间，其强对偶 \(H_{b}^{\prime}\) 具有性质（GDF）（参见附录, No. 2, 命题 3）。令 F 为空间 \(\mathcal{L}_{s}(G;H)\) ；对 T 到 F 中的每个标量本质 \(\mu\)-可积映射 U ，积分 \(\int U d\mu\) 属于 F 。

由于 G 是桶式的（命题 11）， \(\mathcal{L}(\mathrm{G};\mathrm{H}) = \mathcal{L}(\mathrm{G}_{\sigma};\mathrm{H}_{\sigma})\) （TVS, IV, §1, No. 3, 命题 7）；此外，可以把 \(\mathrm{F} = \mathcal{L}_s(\mathrm{G};\mathrm{H})\) 换成空间 \(\mathcal{L}_s(\mathrm{G}_{\sigma};\mathrm{H}_{\sigma})\) ，因为这两个空间有相同的对偶 \(\mathrm{G} \otimes \mathrm{H}'\) （TVS, II, §6, No. 2, 命题 3, 以及上文第 3 目第一段）。若对每个 \(u \in \mathcal{L}(\mathrm{G};\mathrm{H}) = \mathcal{L}(\mathrm{G}_{\sigma};\mathrm{H}_{\sigma})\) ，令 \(\widetilde{u}(\mathbf{x},\mathbf{y}') = \langle u(\mathbf{x},\mathbf{y}') \rangle\) （对 \(\mathbf{x} \in \mathbf{G}\) ， \(y' \in H')\)），则线性映射 \(u \mapsto \widetilde{u}\) 是 F 到空间 \(F_1\) （由 \(G_\sigma \times H_s'\) 上分别连续双线性型组成）上的双射，其中 \(H_s'\) 表示对偶 \(H'\) （赋以弱拓扑 \(\sigma(H', H)\) ）（附录, No. 1）；此外，该映射是 \(\mathcal{L}_s(G_\sigma; H_\sigma)\) 到赋以逐点收敛拓扑的 \(F_1\) 上的同构（出处同上）。但由于按假设 H 是 \(H_b'\) 的对偶， \(F_1\) 也是 \(G \times H_b'\) 上分别连续双线性型所成的空间。于是推论 2 由推论 1 得出。

注意，推论 2 特别适用于 G 为 Banach 空间且 H 为自反 Banach 空间的情形。

\subsection{可测映射与标量可测映射}

若 T 到 Hausdorff 局部凸空间 F 中的映射 f 是标量 \(\mu\)-可测的，一般不能推出 f 是 \(\mu\)-可测的（习题 12）。然而：

命题 12. --- 若 F 是可分的可度量化局部凸空间，则 T 到 F 中的每个标量 \(\mu\)-可测映射 f 都是 \(\mu\)-可测的。

因为， F 可以看作 Banach 空间的可数乘积 \(\prod_{n}E_{n}\) 的子空间（TVS, II, §1, No. 3, 命题 3），并且可以假设 \(\mathrm{pr}_n(\mathbf{F})\) 在 \(\mathbf{E}_n\) 中稠密，后者因此是可分的。对每个 \(n\) ，映射 \(\mathrm{pr}_n \circ \mathbf{f}\) 是标量 \(\mu\)-可测的，因而是 \(\mu\)-可测的（Ch. IV, §5, No. 5, 命题 10 的推论 2），从而 \(\mathbf{f}\) 是 \(\mu\)-可测的（Ch. IV, §5, No. 3, 定理 1）。

命题 13. --- 设 F 是局部凸空间，它是一个由可分、可度量化局部凸空间 \(F_{n}\) 组成的序列的直接极限，且 F 是诸 \(F_{n}\) 之并。令 \(F'\) 为 F 的对偶，赋以拓扑 \(\sigma(\mathbf{F}',\mathbf{F})\) 。则每个标量 \(\mu\)-可测映射 f （ T 到 \(F'\) 中的）都是 \(\mu\)-可测的。

先设 F 是可度量化的且可分的，令 D 为 F 中的可数稠密集。令 \((\mathrm{V}_{n})\) 为 F 中 0 的平衡、凸、开邻域的一个递减基本序列；极集 \(V_{n}^{\circ}\) 是等度连续的，它们的并是整个 \(F'\) 。令 \(\mathrm{T}_{n} = \mathbf{f}^{-1}(\mathrm{V}_{n}^{\circ})\) ；序列 \((\mathrm{T}_{n})\) 是递增的，且 \(T = \bigcup_{n} T_{n}\) ；我们来证明每个 \(T_{n}\) 都是 \(\mu\)-可测的。事实上， \(D \cap V_{n}\) 在 \(V_{n}\) 中稠密；对每个 \(y \in D \cap V_{n}\) ，令 \(S_{y}\) 为 \(t \in T\) 中使 \(|\langle y, f(t) \rangle| \leqslant 1\) 的元素所成之集；假设蕴含每个 \(S_{y}\) 都可测，而 \(T_{n}\) 是诸 \(S_{y}\) \((y \in D \cap V_{n})\) 组成的可数族的交。在此条件下，对 T 的每个紧子集 K 与每个 \(\varepsilon > 0\) ，存在整数 n 使 \(\mu\left(\mathrm{K}-\left(\mathrm{K}\cap\mathrm{T}_{n}\right)\right)\leqslant\frac{\varepsilon}{4}\) ，然后存在紧子集 \(K_{1}\) （含于 \(K\cap T_{n}\) ）使 \(\mu\left(\left(\mathrm{K}\cap\mathrm{T}_{n}\right)-\mathrm{K}_{1}\right)\leqslant\frac{\varepsilon}{4}\) ；最后，存在紧子集 \(K_{2}\) （含于 \(K_{1}\) ）使 \(\mu(K_{1}-K_{2})\leqslant\frac{\varepsilon}{2}\) ，且使在 \(K_{2}\) 上诸函数 \(\langle y,f\rangle\) （ \(y\in D\) ）的限制都连续（Ch. IV, §5, No. 1, 命题 2）。由于集 \(\mathbf{f}(K_{2})\subset\mathbf{f}(T_{n})\subset V_{n}^{\circ}\) 是等度连续的， \(\sigma(F',F)\) 在 \(\mathbf{f}(K_{2})\) 上诱导的拓扑与在 D 中逐点收敛的拓扑相同（GT, X, §2, No. 4, 定理 1）；因此 f 在 \(K_{2}\) 上的限制是连续的，由此得到第一种情形下的断言。

现在过渡到一般情形。若 \(z^{\prime}\) 是 F 上的连续线性形式，其限制 \(z_{n}^{\prime}\) （到 \(F_{n}\) 上）是连续的；由于 \(F = \bigcup_{n} F_{n}\) ，F 的对偶 \(F^{\prime}\) 可以（在代数上）等同于乘积 \(\prod_{n} F_{n}^{\prime}\) 的一个线性子空间，且 \(pr_{n}z^{\prime} = z_{n}^{\prime}\) 。此外，由于 F 的每个有限子集都含于某个 \(F_{n}\) ，拓扑 \(\sigma(F^{\prime}, F)\) 不是别的，正是由诸拓扑 \(\sigma(F_{n}^{\prime}, F_{n})\) 的乘积拓扑所诱导的拓扑。在此条件下，若 f 是标量 \(\mu\)-可测的，则 \(pr_{n} \circ f\) 是标量 \(\mu\)-可测的，因为对每个 \(t \in T\) ， \(\operatorname{pr}_{n}(f(t))\) 是 \(f(t)\) 在 \(F_{n}\) 上的限制。证明的第一部分表明，对每个 n ， \(pr_{n} \circ f\) 是 \(\mu\)-可测的，因此 f 也是（Ch. IV, §5, No. 3, 定理 1）。

\subsection{应用：I. 连续函数到测度空间的延拓}

设 T 为局部紧空间， F 为拟完备的 Hausdorff 局部凸空间， f 为 T 到 F 中的连续映射；若 \(\mu\) 是 T 上的正测度，具有紧支集 S ，则 \(\mathbf{f}(S)\) 是紧的；于是 \(\mathbf{f}(S)\) 的闭凸包络是紧的（TVS, III, §1, No. 6），因此 f 是标量 \(\mu\)-可积的且 \(\int f d\mu \in F\) （No. 2, 命题 5 的推论）。现在若 \(\lambda\) 是任一具有紧支集的实测度，则 \(\lambda^{+}\) 与 \(\lambda^{-}\) 都是具有紧支集的正测度；令 \(\int f d\lambda = \int f d\lambda^{+} - \int f d\lambda^{-}\) ，即可立即验证（利用关系式 \((\lambda + \mu)^{+} + \lambda^{-} + \mu^{-} = \lambda^{+} + \mu^{+} + (\lambda + \mu)^{-}\) ） \(\lambda \mapsto \int f d\lambda\) 是由 T 上具有紧支集的测度所成的空间 \(\mathcal{C}'(T)\) 到局部凸空间 F 中的线性映射。

现在我们注意到，空间 \(\mathcal{C}^{\prime}(\mathrm{T})\) 可以等同于 T 上连续数值函数所成的空间 \(\mathcal{C}(\mathrm{T})\) 的对偶（其记法即由此而来），但须 \(\mathcal{C}(\mathrm{T})\) 赋以紧收敛拓扑（在本目及下一目中我们总作此假设）：因为一方面已知（Ch. IV, §4, No. 8, 命题 14） T 上可以延拓为 \(\mathcal{C}(\mathrm{T})\) 上的连续线性形式的测度恰是具有紧支集的测度；反过来，连续线性形式在 \(\mathcal{K}(T)\) 上的限制（该形式定义在 \(\mathcal{C}(T)\) 上）是一个测度（因为 \(\mathcal{K}(T)\) 的拓扑细于 \(\mathcal{C}(T)\) 的拓扑所诱导的拓扑）。

命题 14. --- 设 T 为局部紧空间， F 为拟完备的 Hausdorff 局部凸空间， f 为 T 到 F 中的连续映射。若 T 上具有紧支集的测度所成的空间 \(\mathcal{C}'(T)\) 赋以在 \(\mathcal{C}(T)\) 的紧子集上一致收敛的拓扑，则映射 \(\lambda \mapsto \int \mathbf{f} d\lambda\) 是唯一的连续线性映射 \(\widetilde{\mathbf{f}}\) （ \(\mathcal{C}'(T)\) 到 F 中），使得 \(\widetilde{\mathbf{f}}(\varepsilon_t) = \mathbf{f}(t)\) 对每个 \(t \in T\) 成立。

为确立延拓的唯一性，只需看出点测度 \(\varepsilon_{t}\) 在 \(\mathcal{C}'(T)\) 中构成一个完全集；由于 \(\mathcal{C}'(T)\) 的对偶是 \(\mathcal{C}(T)\) （TVS, IV, §1, No. 1, 定理 1），只需注意到每个函数 \(g \in \mathcal{C}(T)\) 若与所有测度 \(\varepsilon_{t}\) 正交，则按定义为 0（TVS, IV, §1, No. 2, 命题 2）。

现在证明 \(\lambda \mapsto \int \mathbf{f}d\lambda\) 连续。令 \(V\) 为 \(F\) 中 0 的闭、平衡、凸邻域；只需证明存在相对紧子集 \(L\) （ \(\mathcal{C}(T)\) 中的），使得关系 \(\lambda \in L^{\circ}\) 与 \(\mathbf{z}' \in V^{\circ}\) 蕴含 \(\left|\left\langle \int \mathbf{f}d\lambda, \mathbf{z}' \right\rangle \right| \leqslant 1\) ，或等价地 \(\left|\int \langle \mathbf{f}, \mathbf{z}' \rangle d\lambda \right| \leqslant 1\) 。为此，我们将证明当 \(\mathbf{z}'\) 跑遍 \(V^{\circ}\) 时，集 \(L\) ——由数值函数 \(\langle \mathbf{f}, \mathbf{z}' \rangle\) 组成——在 \(\mathcal{C}(T)\) 中相对紧。由于 \(V^{\circ}\) 对 \(\sigma(F', F)\) 有界，数 \(|\langle \mathbf{f}(t), \mathbf{z}' \rangle|\) （对固定的 \(t \in T\) 与 \(\mathbf{z}'\) 跑遍 \(V^{\circ}\) ）的上确界是有限的；根据 Ascoli 定理（GT, X, §2, No. 5, 定理 2 的推论 2），因此只需证明诸 \(\langle \mathbf{f}, \mathbf{z}' \rangle\) （ \(\mathbf{z}' \in V^{\circ}\) ）所成之集等度连续。但是，对每个 \(t_0 \in T\) 与每个 \(\delta > 0\) ，按假设存在一个邻域 \(W\) （ \(t_0\) 在 \(T\) 中的邻域），使得 \(\mathbf{f}(t) - \mathbf{f}(t_0) \in \delta V\) （对所有 \(t \in W\) ）；由此推出 \(|\langle \mathbf{f}(t), \mathbf{z}' \rangle - \langle \mathbf{f}(t_0), \mathbf{z}' \rangle| \leqslant \delta\) （对所有 \(t \in W\) 与所有 \(\mathbf{z}' \in V^{\circ}\) ），证毕。

注记. --- 1) 映射 \(t \mapsto \varepsilon_t\) 是 T 到空间 \(\mathcal{C}'(\mathrm{T})\) 中的同胚；因为，若 L 是 \(\mathcal{C}(\mathrm{T})\) 的紧子集，且 \(t_0 \in \mathrm{T}\) ，则存在（GT, X, §2, No. 5, 定理 2 的推论 3） \(t_0\) 的一个邻域 W ，使得 \(|g(t) - g(t_0)| \leqslant 1\) （对每个 \(t \in \mathrm{W}\) 与每个函数 \(g \in \mathrm{L}\) ），因此 \(\varepsilon_t - \varepsilon_{t_0} \in \mathrm{L}^\circ\) （对 \(t \in \mathrm{W}\) ），这就证明了 \(t \mapsto \varepsilon_t\) 的连续性（参见 Ch. IV, §4, No. 8, 命题 15）；此外已知逆映射对淡拓扑已经连续（Ch. III, §1, No. 9, 命题 13），从而对在 \(\mathcal{C}(\mathrm{T})\) 的紧子集上一致收敛的拓扑更是如此。如果进而把 T 与它在 \(\mathcal{C}'(\mathrm{T})\) 中的像（借助 \(t \mapsto \varepsilon_t\) ）等同起来，就可以说 \(\lambda \mapsto \int f d\lambda\) 是 f 到线性映射的唯一连续延拓。

\begin{enumerate}
\def\labelenumi{\arabic{enumi})}
\setcounter{enumi}{1}
\tightlist
\item
  注意，在 \(\lambda\mapsto\int f d\lambda\) 的连续性的证明中，我们没有用到 F 拟完备这一事实。因此命题 14 的结论在没有这一假设时仍然成立，只要另外已知 \(\int f d\mu\in F\) 对每个具有紧支集的正测度 \(\mu\) 成立。
\end{enumerate}

现在设 \(\mathbf{f}(\mathrm{T})\) 是 \(\mathrm{F}\) 的有界子集。于是，对每个有界正测度 \(\mu\) （ \(\mathrm{T}\) 上的）， \(\mathbf{f}\) 是标量 \(\mu\)-可积的，且 \(\int \mathbf{f}d\mu \in \mathrm{F}\) （No. 2, 命题 8）。若 \(\lambda\) 是 T 上任一有界实测度，则 \(\lambda^{+}\) 与 \(\lambda^{-}\) 都有界，并且立即可以看出，如上定义的 \(\lambda \mapsto \int \mathbf{f} d\lambda\) 是 T 上有界测度所成的空间 \(\mathcal{M}^1(\mathrm{T})\) 到局部凸空间 F 中的线性映射，它显然延拓了映射 \(\lambda \mapsto \int \mathbf{f} d\lambda\) （ \(\mathcal{C}'(\mathrm{T})\) 到 F 中）。

命题 15. --- 设 T 为局部紧空间， F 为拟完备的 Hausdorff 局部凸空间， f 为 T 到 F 中的连续映射，使得 \(\mathbf{f}(\mathrm{T})\) 有界。若 \(\mathcal{M}^{1}(\mathrm{T})\) 赋以它的 Banach 空间拓扑，则线性映射 \(\lambda \mapsto \int f d\lambda\) （ \(\mathcal{M}^{1}(\mathrm{T})\) 到 F 中）是连续的。

对 F 中 0 的每个闭、平衡、凸邻域 V ，存在 \(\rho > 0\) 使 \(\mathbf{f}(\mathrm{T}) \subset \rho\mathrm{V}\) ；于是 \(\mathbf{f}(\mathrm{T})\) 的闭、平衡、凸包络 B 含于 \(\rho V\) ，且由假设它是完备的。这时若 \(\|\lambda\| \leqslant 1/\rho\) ，则由第 2 目命题 8 及关系式 \(\|\lambda\| = \lambda^{+}(\mathrm{T}) + \lambda^{-}(\mathrm{T})\) 得 \(\int f d\lambda \in B/\rho \subset V\) 。

\subsection{应用：II. 取值于算子空间的连续函数到测度空间的延拓}

设 G 为 Hausdorff 局部凸空间， H 为 Hausdorff 且拟完备的局部凸空间，并用 F 表示由 G 到 H 中的连续线性映射所成的空间 \(\mathcal{L}(\mathrm{G};\mathrm{H})\) ，赋以紧收敛拓扑。空间 F 未必拟完备，且若 \(t \mapsto U(t)\) 是 T 到 F 中的连续映射， \(\mu\) 是 T 上具有紧支集的正测度，则未必有 \(\int U d\mu \in F\) （习题 27）。然而，若对 T 的每个紧子集 K ， \(U(\mathrm{K})\) 都等度连续，则它在 F 中的平衡凸包络也等度连续（TVS, III, §3, No. 4），且由于 H 拟完备，这个平衡凸包络的闭包将是 F 的完备子集（TVS, III, §3, No. 8, 命题 11）；于是确实有 \(\int U d\mu \in F\) （No. 2, 命题 8）。

加在 U 上的补充条件可以用别的方式表述：

引理 3. --- 设 G , H 为两个局部凸空间， U 为局部紧空间 T 到 \(\mathcal{L}(\mathrm{G};\mathrm{H})\) 中的映射。下列条件等价：

\begin{enumerate}
\def\labelenumi{\alph{enumi})}
\item
  映射 \((t,\mathbf{x})\mapsto U(t)\cdot \mathbf{x}\) （由 \(\mathrm{T}\times \mathrm{G}\) 到 H 中）是连续的。
\item
  对 T 的每个紧子集 K ， \(U(\mathbf{K})\) 等度连续，且存在完全集 \(D \subset G\) ，使得对每个 \(x \in D\) ，映射 \(t \mapsto U(t) \cdot x\) 在 T 上连续。
\end{enumerate}

此外，当 \(U\) 满足这些条件时， \(U\) 是 \(\mathrm{T}\) 到赋以紧收敛拓扑的 \(\mathcal{L}(\mathrm{G};\mathrm{H})\) 中的连续映射。

为看出 a) 蕴含 b) ，我们注意到：对 H 中 0 的每个邻域 V 与每个 \(t \in \mathbf{K}\) ，按假设存在邻域 \(\mathrm{L}_t\) （ \(t\) 在 T 中的邻域）与 G 中 0 的邻域 \(\mathrm{W}_t\) ，使得关系 \(t' \in \mathrm{L}_t\) 与 \(\mathbf{x} \in \mathrm{W}_t\) 蕴含 \(U(t') \cdot \mathbf{x} \in \mathrm{V}\) 。只需用有限个邻域 \(\mathrm{L}_{t_i}\) 覆盖 K ，并取 \(\mathrm{W} = \bigcap_{i} \mathrm{W}_{t_i}\) ，就能有 \(U(t) \cdot \mathbf{x} \in \mathrm{V}\) （只要 \(t \in \mathrm{K}\) 与 \(\mathbf{x} \in \mathrm{W}\) ），这证明了 \(U(\mathrm{K})\) 的等度连续性。

反之，设 b) 成立；只需证明对 T 的每个紧子集 K ，映射 \((t,\mathbf{x})\mapsto U(t)\cdot\mathbf{x}\) 在 \(K\times G\) 上连续。令 \(M=U(K)\) ；由于 M 等度连续，在 M 上， G 中逐点收敛拓扑与 D 中逐点收敛拓扑相同（GT, X, §2, No. 4, 定理 1）；因此假设 b) 蕴含 \(t\mapsto U(t)\) 是 K 到 \(\mathcal{L}(G;H)\) 中的连续映射（当 \(\mathcal{L}(G;H)\) 赋以逐点收敛拓扑时）。另一方面， \((A,\mathbf{x})\mapsto A\cdot\mathbf{x}\) 是 \(M\times G\) 到 H 中的连续映射（当 M 赋以逐点收敛拓扑时）（GT, X, §2, No. 1, 命题 1 的推论 4）。由于映射 \((t,\mathbf{x})\mapsto U(t)\cdot\mathbf{x}\) 可以分解为 \((t,\mathbf{x})\mapsto(U(t),\mathbf{x})\mapsto U(t)\cdot\mathbf{x}\) ，我们断定它是连续的。

最后，引理的最后一个断言由如下事实得出：在 M 上，紧收敛拓扑与逐点收敛拓扑相同（GT, X, §2, No. 4, 定理 1）。

于是，设 U 满足引理 3 的条件；那么（若 H 拟完备）可以像第 6 目中那样，定义一个线性映射 \(\lambda \mapsto \int U d\lambda\) （由 \(\mathcal{C}'(\mathrm{T})\) 到 \(\mathrm{F} = \mathcal{L}(\mathrm{G}; \mathrm{H})\) 中）。我们令 \(U(\lambda) = \int U d\lambda\) 。

命题 16. --- 设 G , H 为两个 Hausdorff 局部凸空间，且假设 H 拟完备。设 U 为 T 到 \(\mathcal{L}(\mathrm{G};\mathrm{H})\) 中的映射，使得 \((t,\mathbf{x})\mapsto U(t)\cdot \mathbf{x}\) 是 \(\mathrm{T}\times \mathrm{G}\) 到 H 中的连续映射。则双线性映射 \((\lambda ,\mathbf{x})\mapsto U(\lambda)\cdot \mathbf{x}\) —— \(\mathcal{C}'(\mathrm{T})\times \mathrm{G}\) 到 H 中——关于 \(\mathcal{C}'(\mathrm{T})\) 的等度连续子集与 G 的紧子集是亚连续的（这蕴含线性映射 \(\lambda \mapsto U(\lambda)\) —— \(\mathcal{C}'(\mathrm{T})\) 到 F 中——是连续的）。

\(\lambda \mapsto U(t)\) 作为 \(\mathcal{C}'(\mathrm{T})\) 到 \(\mathrm{F}\) 中的映射的连续性，由引理 3 及第 6 目命题 14 后的注记 2 得出。于是，剩下要证明：对每个闭、平衡、凸邻域 \(\mathrm{V}\) （ \(\mathrm{H}\) 中 0 的）与每个等度连续子集 \(\mathrm{N}\) （ \(\mathcal{C}'(\mathrm{T})\) 的），存在邻域 \(\mathrm{W}\) （ \(\mathrm{G}\) 中 0 的），使得关系 \(\mathbf{x} \in \mathrm{W}\) ， \(\lambda \in \mathrm{N}\) 蕴含 \(U(\lambda) \cdot \mathbf{x} \in \mathrm{V}\) 。可以假设 \(\mathrm{N} = \mathrm{S}^{\circ}\) ，其中 \(\mathrm{S}\) 是 \(\mathcal{C}(\mathrm{T})\) 中 0 的邻域，从而可以假设 \(\mathrm{S}\) 是由满足如下条件的函数 \(g \in \mathcal{C}(\mathrm{T})\) 所成的集：\(|g(t)| \leqslant 1\) 在某个紧子集 \(\mathrm{K}\) （ \(\mathrm{T}\) 的）上成立。只需证明 \(|\langle U(\lambda) \cdot \mathbf{x}, \mathbf{x}' \rangle| \leqslant 1\) （对 \(\mathbf{x} \in W\) ， \(\mathbf{x}' \in V^\circ\) 与 \(\lambda \in S^\circ\) ）。而今，由于 \(U(K)\) 等度连续，存在邻域 \(W\) （ \(G\) 中 0 的），使得关系 \(t \in K\) ， \(\mathbf{x} \in W\) 蕴含 \(U(t) \cdot \mathbf{x} \in V\) ；因此关系 \(\mathbf{x} \in W\) ， \(\mathbf{x}' \in V^\circ\) 蕴含函数 \(t \mapsto \langle U(t) \cdot \mathbf{x}, \mathbf{x}' \rangle\) 属于 \(S\) ，从而 \(|\langle U(t) \cdot \mathbf{x}, \mathbf{x}' \rangle| = |\int\langle U(t) \cdot \mathbf{x}, \mathbf{x}' \rangle d\lambda(t)| \leqslant 1\) （按 \(S^\circ\) 的定义）。

现在假设 \(U\) 是 \(\mathrm{T}\) 到 \(\mathrm{F}\) 中的连续映射，而且 \(U(\mathrm{T})\) 等度连续。于是，与上面相同的论证表明（由于 \(\mathrm{H}\) 拟完备），对每个有界正测度 \(\mu\) （ \(\mathrm{T}\) 上的）， \(\int U d\mu \in \mathrm{F}\) 。因此可以像上面那样定义线性映射 \(\lambda \mapsto \int U d\lambda = U(\lambda)\) （ \(\mathcal{M}^1(\mathrm{T})\) 到 \(\mathrm{F}\) 中），它延拓了 \(\mathcal{C}'(\mathrm{T})\) 到 \(\mathrm{F}\) 中的类似映射。此外，对每个闭、平衡、凸邻域 \(\mathrm{V}\) （ \(\mathrm{H}\) 中 0 的），按假设存在邻域 \(\mathrm{W}\) （ \(\mathrm{G}\) 中 0 的），使得对每个 \(\mathbf{x} \in \mathrm{W}\) 与每个 \(t \in \mathrm{T}\) ，有 \(U(t) \cdot \mathbf{x} \in \mathrm{V}\) ，从而（由于 \(\mathrm{V}\) 弱闭） \(\int (U(t) \cdot \mathbf{x}) d\lambda(t) \in \| \lambda \| \cdot \mathrm{V}\) （No. 2, 命题 5）。换言之：

命题 17. --- 设 G , H 为两个 Hausdorff 局部凸空间，且假设 H 拟完备。设 U 为 T 到 \(\mathcal{L}(G;H)\) 中的映射，使得 \((t,\mathbf{x})\mapsto U(t)\cdot \mathbf{x}\) 在 T×G 上连续，且 \(U(T)\) 等度连续。则当 \(\mathcal{M}^1 (T)\) 赋以它的 Banach 空间拓扑时，双线性映射 \((\lambda ,\mathbf{x})\mapsto U(\lambda)\cdot \mathbf{x}\) —— \(\mathcal{M}^1 (T)\times G\) 到 H 中——是连续的（特别地，这蕴含线性映射 \(\lambda \mapsto U(\lambda)\) —— \(\mathcal{M}^1 (T)\) 到 \(\mathcal{L}(G;H)\) 中——当 \(\mathcal{L}(G;H)\) 赋以有界收敛拓扑时是连续的）。

命题 18. --- 设 \(G_{1}, G_{2}, H_{1}, H_{2}\) 为四个 Hausdorff 局部凸空间，其中 \(H_{1}\) 与 \(H_{2}\) 假设为拟完备的。令 \(A : G_{1} \to G_{2}\) 与 \(B : H_{1} \to H_{2}\) 为两个连续线性映射。令 \(U_{1} : T \to \mathcal{L}(G_{1}; H_{1})\) ， \(U_{2} : T \to \mathcal{L}(G_{2}; H_{2})\) 为两个满足命题 16（相应地，命题 17）条件的映射，并假设对每个 \(t \in T\) ， \(B \circ U_{1}(t) = U_{2}(t) \circ A\) 。则对 T 上每个具有紧支集的（相应地，有界的）测度 \(\lambda\) ， \(B \circ U_{1}(\lambda) = U_{2}(\lambda) \circ A\) 。

事实上，对每个 \(x \in G_{1}\) ，有（No. 1, 命题 1）

\[
\begin{array}{c} \left(B \circ U _ {1} (\lambda)\right) \cdot \mathbf {x} = \int \Big (\left(B \circ U _ {1} (t)\right) \cdot \mathbf {x} \Big) d \lambda (t) \\ = \int \Big (\left(U _ {2} (t) \circ A\right) \cdot \mathbf {x} \Big) d \lambda (t) = U _ {2} (\lambda) \cdot (A \cdot \mathbf {x}). \end{array}
\]

注记. --- 1) 设 G 与 H 为 Banach 空间， U 为 T 到 \(\mathcal{L}(\mathrm{G};\mathrm{H})\) 中的映射，使得 \((t,\mathbf{x})\mapsto U(t)\cdot\mathbf{x}\) 在 \(T\times G\) 上连续。注意，这蕴含有限函数 \(t\mapsto\|U(t)\|\) 在 T 的每个紧子集上有界且在 T 上下半连续，因为它是连续函数 \(t \mapsto |U(t) \cdot \mathbf{x}|\) （x 跑遍球 \(|x| \leqslant 1\) ， G 中的）的上包络。令 \(h(t) = \|U(t)\|\) 。于是，对每个使 h 为 \(\mu\)-可积的正测度 \(\mu\) （ T 上），我们仍有 \(\int U d\mu \in \mathcal{L}(\mathrm{G};\mathrm{H})\) 。因为，测度 \(\nu = h \cdot \mu\) 按假设有界；因此存在 T 的一个分割，由一个 \(\nu\)-可忽略集 N 与紧子集序列 \((\mathrm{K}_{n})\) 组成。把本目开头所作的论证应用于测度 \(\varphi_{K_{n}} \cdot \mu\) ，便表明

\[
A _ {n} = \int \varphi_ {\mathrm{K} _ {n}} U d \mu \in \mathrm{F} = \mathcal {L} (\mathrm{G}; \mathrm{H}),
\]

而且（No. 2, 命题 6）， \(\| A_n\| \leqslant \int \varphi_{\mathrm{K}_n}\| U\| d\mu \leqslant \nu (\mathrm{K}_n)\) 。因此以 \(A_{n}\) 为通项的级数在 Banach 空间 \(\mathcal{L}(\mathrm{G};\mathrm{H})\) 中绝对收敛，并且立即可以看出其和为 \(\int Ud\mu\) ，且 \(\left\| \int Ud\mu \right\| \leqslant \int \| U\| d\mu\) 。

\begin{enumerate}
\def\labelenumi{\arabic{enumi})}
\setcounter{enumi}{1}
\tightlist
\item
  设 \(\mathrm{G} = \mathrm{H}\) 拟完备，且 \(U\) 满足命题 16 的假设。令 \(\mathbf{M}\) 为空间 \(\mathcal{C}'(\mathrm{T})\) 的一个稠密子集（对弱拓扑 \(\sigma\big(\mathcal{C}'(\mathrm{T}),\mathcal{C}(\mathrm{T})\big)\) 而言），令 \(\mathrm{X}\) 为 \(\mathrm{H}\) 的闭线性子空间，使得 \(U(\lambda)(\mathrm{X})\subset \mathrm{X}\) （对每个测度 \(\lambda \in \mathrm{M}\) ）。则也有 \(U(t)(\mathrm{X})\subset \mathrm{X}\) （对每个 \(t\in \mathrm{T}\) ）：事实上，对每个 \(\mathbf{x}\in \mathrm{X}\) 与每个 \(\mathbf{x}'\in \mathrm{H}'\) （与 \(\mathrm{X}\) 正交），按假设 \(\langle U(\lambda)\cdot \mathbf{x},\mathbf{x}'\rangle = 0\) （对所有 \(\lambda \in \mathrm{M}\) ），它可以写成 \(\int \langle U(t)\cdot \mathbf{x},\mathbf{x}'\rangle d\mu (t) = 0\) 。连续函数 \(t\mapsto \langle U(t)\cdot \mathbf{x},\mathbf{x}'\rangle\) 由于与 \(\mathrm{M}\) 正交而为 0 ，由此得到 \(\langle U(t)\cdot \mathbf{x},\mathbf{x}'\rangle = 0\) （对每个 \(\mathbf{x}'\in \mathrm{X}^\circ\) ），从而 \(U(t)\cdot \mathbf{x}\in \mathrm{X}\) （对所有 \(t\in \mathrm{T}\) 与 \(\mathbf{x}\in \mathrm{X}\) ），这就证明了我们的断言。
\end{enumerate}

\section{向量测度}

\subsection{向量测度的定义}

Ch. III, §1, No. 3 中给出的测度定义可以如下推广：

定义 1. --- 设 F 是 R 上的 Hausdorff 局部凸空间。把空间 \(\mathcal{H}(T)\) 到 F 中的每个连续线性映射都称为 T 上取值于 F 的向量测度。

定义 1 也可以表述如下：T 上取值于 F 的向量测度是一个线性映射 m （由 \(\mathcal{K}(T)\) 到 F 中），使得对 T 的每个紧子集 K ， m 在 \(\mathcal{K}(T,K)\) 上的限制对一致收敛拓扑连续。若 \(f \in \mathcal{K}(T)\) ，也用 \(\int f dm\) 或 \(\int f(t) dm(t)\) 代替 \(\mathbf{m}(f)\) 来记。取值于 R 中的测度有时称为实测度（Ch. III, §1, No. 5）或 T 上的标量测度。

例子. --- 1) \(\mathcal{H}(\mathrm{T})\) 的恒等映射是 \(\mathrm{T}\) 上取值于 \(\mathcal{H}(\mathrm{T})\) 的向量测度。

\begin{enumerate}
\def\labelenumi{\arabic{enumi})}
\setcounter{enumi}{1}
\tightlist
\item
  *设 H 是复 Hilbert 空间， \(\mathcal{L}(\mathrm{H})\) 是 H 的连续自同态所成的赋范代数。设 A 是 \(\mathcal{L}(\mathrm{H})\) 的一个子代数，它交换、闭、自伴且包含单位元；可以证明，存在一个紧空间 X 以及赋范代数 A 到代数 \(\mathcal{K}_{\mathbf{C}}(\mathrm{X}) = \mathcal{C}_{\mathbf{C}}(\mathrm{X})\) （ X 上连续复函数所成的代数）上的一个同构，该代数赋以范数 \(\| f \| = \sup_{x \in \mathrm{X}} |f(x)|\) 。逆同构在 \(\mathcal{K}(\mathrm{X})\) 上的限制是一个向量测度 \(\mathbf{m}\) （ X 上的，取值于 \(\mathcal{L}(\mathrm{H})\) ），使得 \(\mathbf{m}(fg) = \mathbf{m}(f)\mathbf{m}(g)\) 。*
\end{enumerate}

注记. --- 1) 为使线性映射 m （由 \(\mathcal{K}(\mathrm{T})\) 到 F 中）是向量测度，必须且只需：对 T 的每个紧子集 K ， m 下单位球 \(\|f\|\leqslant1\) （ \(\mathcal{K}(\mathrm{T},\mathrm{K})\) 的）的像在 F 中有界。因此，取值于 F 的向量测度这一概念，对 F 上具有相同有界集的所有 Hausdorff 局部凸拓扑都是一样的，特别地，对与 F 和 F' 之间的对偶相容的所有拓扑都是一样的（TVS, IV, §1, No. 1, 命题 1）。

\begin{enumerate}
\def\labelenumi{\arabic{enumi})}
\setcounter{enumi}{1}
\item
  设 \(T_{1}\) 为局部紧空间， \(F_{1}\) 为 R 上的 Hausdorff 局部凸空间， u 为 \(\mathcal{K}(T_{1})\) 到 \(\mathcal{K}(T)\) 中的连续线性映射， v 为 F 到 \(F_{1}\) 中的连续线性映射。若 m 是 T 上取值于 F 的向量测度，则 \(v \circ m \circ u\) 是 \(T_{1}\) 上取值于 \(F_{1}\) 的向量测度。特别地，若 g 是 T 上的连续有限数值函数，则 \(f \mapsto \mathbf{m}(gf)\) 是 T 上取值于 F 的向量测度，记作 \(g \cdot m\) ；若 h 是 T 上另一个连续有限数值函数，则 \(g \cdot (h \cdot \mathbf{m}) = (gh) \cdot \mathbf{m}\) 。
\item
  由于空间 \(\mathcal{H}(\mathrm{T})\) 是诸 Banach 空间 \(\mathcal{H}(\mathrm{T},\mathrm{K})\) 的直接极限，因而特别是桶式的（TVS, III, §4, No. 1, 命题 2 的推论及命题 3 的推论 3），为使线性映射 \(\mathbf{m}\) （由 \(\mathcal{H}(\mathrm{T})\) 到 \(\mathbf{F}\) 中）是向量测度，必须且只需：对每个 \(\mathbf{z}' \in \mathbf{F}'\) ， \(\mathbf{z}' \circ \mathbf{m}\) 是 \(\mathrm{T}\) 上的标量测度（TVS, II, §6, No. 4, 命题 5 以及 IV, §1, No. 3, 命题 7）。
\item
  鉴于注记 1，Ch. III, §2, No. 1 的命题 1 及其证明对向量测度仍然有效。因此可以把向量测度 \(\mathbf{m}\) （ \(T\) 上的）的支集定义为最大的开集 \(U \subset T\) （使 \(\mathbf{m}\) 在 \(U\) 上的限制为零者）的余集。
\end{enumerate}

\subsection{关于向量测度的积分}

设 \(\mathbf{m}\) 是 \(T\) 上取值于 \(F\) 的向量测度。对每个 \(\mathbf{z}' \in F'\) ，映射 \(\mathbf{z}' \circ \mathbf{m}\) 是 \(T\) 上的标量测度，且线性地依赖于 \(\mathbf{z}'\) 。若 \(f\) 是定义在 \(T\) 上的数值函数，我们约定（滥用语言）说偶 \((f, \mathbf{m})\) 具有性质 \(\boldsymbol{P}\) ，如果对每个 \(\mathbf{z}' \in F'\) ，偶 \((f, |\mathbf{z}' \circ \mathbf{m}|)\) 具有性质 \(\boldsymbol{P}\) 。例如，我们说 \(f\) 关于 \(\mathbf{m}\) 本质可积，意思是对每个 \(\mathbf{z}' \in F'\) ，函数 \(f\) 关于 \(|\mathbf{z}' \circ \mathbf{m}|\) 本质可积。这等价于说 \(f\) 关于测度 \((\mathbf{z}' \circ \mathbf{m})^+\) 与 \((\mathbf{z}' \circ \mathbf{m})^-\) 中的每一个都本质可积（Ch. V, §2, No. 2, 命题 3 的推论 2）。

命题 1. --- 设 m 是 T 上取值于 F 的向量测度， f 是 T 上的数值函数，关于 m 本质可积。映射

\[
\mathbf {z} ^ {\prime} \mapsto \int f d (\mathbf {z} ^ {\prime} \circ \mathbf {m}) ^ {+} - \int f d (\mathbf {z} ^ {\prime} \circ \mathbf {m}) ^ {-}
\]

是 \(\mathbf{F}'\) 上的线性形式。

用 \(\Phi\) 记这一映射，立即有 \(\Phi(\lambda\mathbf{z}') = \lambda\Phi(\mathbf{z}')\) （对一切 \(\lambda \in \mathbf{R}\) ）。一切归结为证明 \(\Phi(\mathbf{y}' + \mathbf{z}') = \Phi(\mathbf{y}') + \Phi(\mathbf{z}')\) 。令 \(\mu = |\mathbf{y}' \circ \mathbf{m}| + |\mathbf{z}' \circ \mathbf{m}|\) ；由 Lebesgue-Nikodym 定理，这时可以写 \(\mathbf{y}' \circ \mathbf{m} = g \cdot \mu\) 与 \(\mathbf{z}' \circ \mathbf{m} = h \cdot \mu\) ，其中 \(g\) 与 \(h\) 是两个有界且局部 \(\mu\)-可积的数值函数（Ch. V, §5, No. 5, 定理 2）；此外， \((\mathbf{y}' \circ \mathbf{m})^+ = g^+ \cdot \mu\) 与 \((\mathbf{y}' \circ \mathbf{m})^- = g^- \cdot \mu\) ，并且当把 \(\mathbf{y}'\) 换成 \(\mathbf{z}'\) （相应地， \(\mathbf{y}' + \mathbf{z}'\) ）、把 \(g\) 换成 \(h\) （相应地， \(g + h\) ）时，类似的关系式成立。在此条件下，立即看出 \(f\) 本质 \(\mu\)-可积（Ch. V, §2, No. 2, 命题 3 的推论 1），而要证明的关系式归结为 \((g + h)^+ - (g + h)^- = (g^+ - g^-) + (h^+ - h^-)\) ，这是显然的。

定义 2. --- 设 m 是 T 上取值于 F 的向量测度， f 是 T 上的数值函数，关于 m 本质可积。把 f 关于 m 的积分（记作 \(\mathbf{m}(f)\) 或 \(\int f dm\) 或 \(\int f(t) dm(t)\) ）定义为 \(F'^{*}\) 中由下式给出的元素：

\[
\left\langle \mathbf {z} ^ {\prime}, \int f d \mathbf {m} \right\rangle = \int f d (\mathbf {z} ^ {\prime} \circ \mathbf {m}) ^ {+} - \int f d (\mathbf {z} ^ {\prime} \circ \mathbf {m}) ^ {-}.\tag{1}
\]

我们注意到，若 \(f \in \mathcal{K}(T)\) ，则如此定义的元素 \(\int f dm\) 与第 1 目中同样记号的元素一致，因为这时 (1) 的右端按定义就是 \(\int f d(\mathbf{z}' \circ \mathbf{m}) = \mathbf{z}'(\mathbf{m}(f))\) 。此外，若特别把定义 2 应用于 F = R 的情形，则可以看出，对每个 \(z' \in F'\) ， \(f\) 关于实测度 \(\mathbf{z}' \circ \mathbf{m}\) 本质可积，且 (1) 的右端可以写成 \(\int f d(\mathbf{z}' \circ \mathbf{m})\) 。

现在设 \(f\) 关于 \(\mathbf{m}\) 本质可积，并令 \(\mathbf{z}' \in \mathbf{F}'\) 。令 \(\mu = |\mathbf{z}' \circ \mathbf{m}|\) ；由 Lebesgue-Nikodym 定理，可以写 \(\mathbf{z}' \circ \mathbf{m} = g \cdot \mu\) ，其中 \(g\) 局部 \(\mu\)-可积且 \(\| g \| \leqslant 1\) ，而命题 1 的证明表明 \(\int f d(\mathbf{z}' \circ \mathbf{m}) = \int fg d\mu\) 。由此，

\[
\left| \int f d (\mathbf {z} ^ {\prime} \circ \mathbf {m}) \right| \leqslant \int | f | d | \mathbf {z} ^ {\prime} \circ \mathbf {m} |.\tag{2}
\]

显然，关于 m 本质可积的有限数值函数所成之集是 R 上的向量空间；我们把这个空间记作 \(\mathcal{L}(\mathbf{m})\) ，赋以使所有线性形式 \(f \mapsto \int f d(\mathbf{z}' \circ \mathbf{m})\) （其中 \(z'\) 跑遍 \(F'\) ）都连续的最粗局部凸拓扑。注意，一般说来局部凸空间 \(\mathcal{L}(\mathbf{m})\) 不是 Hausdorff 的。

例. --- 取 m 为 \(\mathcal{H}(\mathrm{T})\) 到自身的恒等映射。由于 \(\mathcal{H}(\mathrm{T})\) 的对偶是 T 上标量测度所成的空间 \(\mathcal{M}(\mathrm{T})\) ，故函数 \(f \in \mathcal{L}(\mathbf{m})\) 就是那些关于每个标量测度 \(\mu\) 都本质可积的函数（参见习题 1），而积分 \(\int f dm\) 是线性形式 \(\mu \mapsto \int f d\mu\) （ \(\mathcal{M}(\mathrm{T})\) 上的）。关系式 \(\int f d\mu = 0\) （对每个测度 \(\mu \in \mathcal{M}(\mathrm{T})\) ）仅当 \(f = 0\) 时成立，取 \(\mu = \varepsilon_t\) （其中 \(t\) 是 T 中任意的）即可看出；换言之，映射 \(f \mapsto \int f dm\) 是 \(\mathcal{L}(\mathbf{m})\) 到 \(\mathcal{M}(\mathrm{T})\) 的代数对偶中的单射，它延拓了 \(\mathcal{H}(\mathrm{T})\) 的恒等映射。因此，关系式 \(\int f dm \in \mathrm{F} = \mathcal{H}(\mathrm{T})\) 等价于 \(f \in \mathcal{H}(\mathrm{T})\) 。

设 \(u\) 是 \(F\) 到 Hausdorff 局部凸空间 \(G\) 中的连续线性映射，我们仍用 \(u\) 表示它的经双重转置的延拓（ \(F'^*\) 到 \(G'^*\) 中的线性映射）（§1, No. 1）。在此约定下：

命题 2. --- 每个数值函数 \(f\) （关于 \(\mathbf{m}\) 本质可积的）都关于 \(u \circ \mathbf{m}\) 本质可积，且 \(\int f d(u \circ \mathbf{m}) = u\left(\int f dm\right)\) 。

鉴于等式 \(y'\circ u\circ m = {}^{t}u(y')\circ m\) （对一切 \(y \in G'\) ），命题是显然的。

一般地，若 \(f \in \mathcal{L}(\mathbf{m})\) ，积分 \(\int f dm\) 属于 \(\mathbf{F}'*\) 而不属于 \(\mathbf{F}\) （见上面的例）。然而：

命题 3. --- 若由 \(f \in \mathcal{K}(\mathbf{T})\) （满足 \(\sup_{t \in \mathbf{T}} |f(t)| \leqslant 1\) 者）所成之集在 m 下的像在 F 中相对弱紧，则 \(\int f dm \in \mathbf{F}\) （对每个有界的、关于 m 本质可积的数值函数 \(f\) ）。

令 \(\mathbf{A}\) 为由 \(f\in \mathcal{L}(\mathbf{m})\) （满足 \(\sup_{t\in \mathrm{T}}|f(t)|\leqslant 1\) 者）所成之集，并令

\(B = A \cap \mathcal{K}(T)\) ；按假设， \(\mathbf{m}(B)\) 在 \(F\) 中相对弱紧，因此只需证明 \(\mathbf{m}(A)\) 含于在 \(F'^*\) 中取的 \(\mathbf{m}(B)\) 的闭包（对拓扑 \(\sigma(F'^*, F')\) ）；由于 \(\mathbf{m}(B)\) 是凸的且平衡的，只需证明 \(\mathbf{m}(\mathrm{B})\) 在 \(\mathrm{F}'\) 中的极集含于 \(\mathbf{m}(\mathrm{A})\) 的极集（TVS, II, §6, No. 3, 定理 1）。而今，为使线性形式 \(\mathbf{z}' \in \mathrm{F}'\) 属于 \((\mathbf{m}(\mathrm{B}))^\circ\) ，必须且只需 \(|\langle \mathbf{z}', \mathbf{m}(g) \rangle| = |\int g d(\mathbf{z}' \circ \mathbf{m})| \leqslant 1\) （对每个函数 \(g \in \mathrm{B}\) ），这就意味着标量测度 \(|\mathbf{z}' \circ \mathbf{m}|\) 有界且范数 \(\leqslant 1\) （Ch. III, §1, No. 8）；但由 (2)，后一条件蕴含 \(|\langle \mathbf{z}', \mathbf{m}(f) \rangle| \leqslant 1\) （对每个函数 \(f \in \mathrm{A}\) ），从而 \(\mathbf{z}' \in (\mathbf{m}(\mathrm{A}))^\circ\) 。

推论 1. --- 若对 T 的每个紧子集 K ，由 \(f \in \mathcal{K}(\mathrm{T},\mathrm{K})\) （满足 \(\sup_{t \in \mathrm{T}} |f(t)| \leqslant 1\) 者）所成之集在 m 下的像在 F 中相对弱紧，则 \(\int f dm \in \mathrm{F}\) （对每个有界且具有紧支集的函数 \(f \in \mathcal{L}(\mathbf{m})\) ），且 \(\int f dm \in \mathrm{F}''\) （对每个函数 \(f \in \mathcal{L}(\mathbf{m})\) ）。

第一个断言可由命题 3 立即得出：若 \(f\) 有界且具有紧支集，且 \(U\) 是 \(f\) 的支集的一个相对紧开邻域，则 \(\mathbf{m}\) 在子空间 \(\mathcal{K}(U)\) 上的限制是一个测度 \(\mathbf{m}_U\) （ \(U\) 上的），它满足命题 3 的条件，且 \(\int f dm_U = \int f dm\) （Ch. V, §7, No. 1, 定理 1），因此 \(\int f dm \in F\) 。

现在设 \(f\) 是 \(\mathcal{L}(\mathbf{m})\) 的任意元素；对每个紧子集 \(K\) （ \(T\) 的）与每个整数 \(n > 0\) ，令 \(f_{n,K}\) 为 \(T\) 上如下定义的数值函数：若 \(t \notin K\) ， \(f_{n,K}(t) = 0\) ；若 \(t \in K\) 且 \(|f(t) \leqslant n\) ， \(f_{n,K}(t) = f(t)\) ；最后，若 \(t \in K\) 且 \(|f(t)| > n\) ， \(f_{n,K}(t) = nf(t)/|f(t)|\) 。显然，对每个 \(t \in T\) ， \(f(t)\) 是 \(f_{n,K}(t)\) 关于 Fréchet 滤子与 \(T\) 的紧子集（递增有向）序集的截口滤子的乘积滤子的极限。由于 \(|f_{n,K}| \leqslant |f|\) ，把 Lebesgue 定理与 Ch. V, §1, No. 3 命题 10 应用于每个标量测度 \(|\mathbf{z}' \circ \mathbf{m}|\) ，便知 \(f_{n,K}\) 收敛于 \(f\) （在 \(\mathcal{L}(\mathbf{m})\) 中，关于上述滤子）。因此，积分 \(\int f dm\) 属于在 \(F'^*\) 中（对拓扑 \(\sigma(F'^*, F')\) ）取的集 \(M\) ——由诸 \(m(f_{n,K})\) 组成——的闭包。但推论的第一部分表明 \(M \subset F\) ，而且，对每个 \(z' \in F'\) 有 \(|\langle z', m(f_{n,K})\rangle| \leqslant \int |f| d|\mathbf{z}' \circ \mathbf{m}|\) ，这表明 \(M\) 在 \(F_\sigma\) 中有界，从而在 \(F\) 中也有界（TVS, IV, §1, No. 1, 命题 1）。于是 §1 第 2 目的引理 1 表明 \(\int f dm \in F''\) 。

推论 2. --- 若 F 是半自反的，则 \(\int f dm \in F\) （对每个数值函数 \(f\) ——关于 \(\mathbf{m}\) 本质可积的）。

\subsection{可优超的向量测度}

设 \(q\) 是 \(F\) 上的下半连续半范数。我们用 \(A_q'\) 表示由 \(z' \in F'\) （满足 \(|\langle z', x \rangle| \leqslant q(x)\) ，对一切 \(x \in F\) 者）所成之集。它是在 \(\mathbf{F}'\) 中由 \(\mathbf{x} \in \mathbf{F}\) （满足 \(q(\mathbf{x}) \leqslant 1\) 者）所成之集的极集；对每个 \(\mathbf{x} \in \mathbf{F}\) ， \(q(\mathbf{x}) = \sup_{\mathbf{z}' \in \hat{\mathrm{A}}_q'} |\langle \mathbf{x}, \mathbf{z}' \rangle|\) 。

定义 3. --- 设 m 是 T 上取值于 F 的向量测度。若 q 是 F 上的下半连续半范数，则称 m 是 q-可优超的，如果存在正测度 \(\mu\) ，使得 \(|z^{\prime} \circ m| \leqslant \mu\) （对每个 \(z^{\prime} \in A_{q}^{\prime}\) ）；这时把诸测度 \(|z^{\prime} \circ m|\) （ \(z^{\prime}\) 跑遍 \(A_{q}^{\prime}\) ）的上确界（Ch. III, §2, No. 4, 定理 3）记作 \(q(\mathbf{m})\) 。若 m 关于 F 上的每个连续半范数 q 都是 q-可优超的，则称 m 是可优超的。

若 \(\mathbf{m}\) 与 \(\mathbf{m}'\) 都是 \(q\)-可优超的，则立即可以看出 \(\mathbf{m} + \mathbf{m}'\) 也是 \(q\)-可优超的，并且

\[
q (\mathbf {m} + \mathbf {m} ^ {\prime}) \leqslant q (\mathbf {m}) + q (\mathbf {m} ^ {\prime}).
\]

当 F 是赋范空间、范数记作 \(|x|\) 时，说 m 可优超，就是说诸测度 \(|z'\circ m|\) （其中 \(|z'| \leqslant 1\) ）有同一个正测度作为上界\(^{1}\)；这时就把该测度族的上确界记作 \(|m|\) 。

若 \(\mathbf{F} = \mathbf{R}\) ，则测度 \(|m|\) （对应于欧几里得范数 \(|x|\) ， \(\mathbf{R}\) 上的）与 Ch. III, §1, No. 5 中记作 \(|m|\) 的测度一致。

命题 4. --- 设 \((\mathrm{F}_i)_{1 \leqslant i \leqslant n}\) 是有限族 Hausdorff 局部凸空间， \(\mathbf{F} = \prod_{i=1}^{n} \mathbf{F}_i\) 是它们的乘积， \(q_i\) ( \(1 \leqslant i \leqslant n\) ) 是 \(\mathbf{F}_i\) 上的下半连续半范数，而 \(q\) 是 \(\mathbf{F}\) 上由 \(q(\mathbf{x}_i, \ldots, \mathbf{x}_n) = \sum_{i=1}^{n} q_i(\mathbf{x}_i)\) 定义的半范数。若 \(\mathbf{m}_i\) ( \(1 \leqslant i \leqslant n\) ) 是 \(\mathbf{T}\) 上取值于 \(\mathbf{F}_i\) 的、 \(q_i\)-可优超的向量测度，则测度 \(\mathbf{m} = (\mathbf{m}_1, \ldots, \mathbf{m}_n)\) （取值于 \(\mathbf{F}\) ）是 \(q\)-可优超的。

因为，对偶 \(\mathbf{F}'\) 可以等同于 \(\prod_{i=1}^{n} \mathrm{F}_i'\) ，使得若 \(\mathbf{x} = (\mathbf{x}_i) \in \mathrm{F}\) ， \(\mathbf{z}' = (\mathbf{z}_i') \in \mathrm{F}'\) ，则 \(\langle \mathbf{x}, \mathbf{z}' \rangle = \sum_{i=1}^{n} \langle \mathbf{x}_i, \mathbf{z}_i' \rangle\) 。若 \(|\langle \mathbf{x}, \mathbf{z}' \rangle| \leqslant q(\mathbf{x})\) （对每个 \(\mathbf{x} \in \mathrm{F}\) ），则特别地 \(|\langle \mathbf{x}_i, \mathbf{z}_i' \rangle| \leqslant q_i(\mathbf{x}_i)\) （对 \(1 \leqslant i \leqslant n\) ），而其逆是显然的，这表明集 \(\mathrm{A}_q'\) 是诸 \(\mathrm{A}_{q_i}'\) 的乘积。由于按假设 \(|\mathbf{z}_i' \circ \mathbf{m}_i| \leqslant q_i(\mathbf{m}_i)\) （对 \(\mathbf{z}_i' \in \mathrm{A}_{q_i}'\) ），由此得出

\[
\left| \mathbf {z} ^ {\prime} \circ \mathbf {m} \right| \leqslant \sum_ {i = 1} ^ {n} \left| \mathbf {z} _ {i} ^ {\prime} \circ \mathbf {m} _ {i} \right| \leqslant \sum_ {i = 1} ^ {n} q _ {i} (\mathbf {m} _ {i})
\]

（对每个 \(\mathbf{z}' \in \mathrm{A}_q'\) ），这就证明了命题。

推论. --- 若空间 F 是有限维的，则每个取值于 F 的向量测度 m 都是可优超的。为使一个数值函数关于 m 本质可积，必须且只需它关于 \(|m|\) 本质可积（其中 \(|x|\) 表示 F 上的任一范数）。

命题 5. --- 设 q 是 F 上的下半连续半范数。设 m 是 q-可优超的测度， f 是关于 m 本质可积且使 \(\int f dm \in F\) 的函数。则

\[
q \left(\int f d \mathbf {m}\right) \leqslant \int^ {\bullet} | f | d q (\mathbf {m}).
\]

因为，

\[
q \left(\int f d \mathbf {m}\right) = \sup _ {\mathbf {z} ^ {\prime} \in \mathrm{A} _ {q} ^ {\prime}} \left| \left\langle \mathbf {z} ^ {\prime}, \int f d \mathbf {m} \right\rangle \right| \leqslant \sup _ {\mathbf {z} ^ {\prime} \in \mathrm{A} _ {q} ^ {\prime}} \int | f | d | \mathbf {z} ^ {\prime} \circ m | \leqslant \int^ {\bullet} | f | d q (\mathbf {m})
\]

这是根据 (1) 以及关系式 \(|\mathbf{z}' \circ \mathbf{m}| \leqslant q(\mathbf{m})\) （对 \(\mathbf{z}' \in \mathrm{A}_q'\) ）。

命题 6. --- 设 F 是拟完备的 Hausdorff 局部凸空间， m 是 T 上取值于 F 的可优超测度。若数值函数 f 关于所有测度 \(q(\mathbf{m})\) （其中 q 跑遍 F 上连续半范数所成之集）都本质可积，则 f 关于 m 本质可积，且 \(\int f dm \in F\) 。

我们将利用下面的辅助结果。设 \((\mu_{\iota})_{\iota\in\mathrm{I}}\) 是 T 上正测度的递增有向族。用 \(\mathcal{L}^{1}\big((\mu_{\iota})_{\iota\in\mathrm{I}}\big)\) 表示 T 上有限数值函数中本质 \(\mu_{\iota}\)-可积（对每个 \(\iota\in I\) ）者所成的向量空间，赋以由半范数 \(f\mapsto\mu_{\iota}(|f|)\;\ (\iota\in\mathrm{I})\) 定义的拓扑。设 \(L_{0}\) 是 \(\mathcal{L}^{1}\big((\mu_{\iota})_{\iota\in\mathrm{I}}\big)\) 中由诸乘积 \(g\varphi_{K}\) 生成的线性子空间，其中 g 跑遍 T 上连续有限数值函数所成之集， K 跑遍 T 的紧子集所成之集。

引理 1. --- 当 \(\mathcal{L}_0\) 与 \(\mathcal{K}(\mathrm{T})\) 赋以 \(\mathcal{L}^1\big((\mu_\iota)_{\iota \in \mathrm{I}}\big)\) 的拓扑诱导的拓扑时：

\begin{enumerate}
\def\labelenumi{\alph{enumi})}
\item
  \(\mathcal{L}_0\) 的每个元素都属于 \(\mathcal{K}(\mathrm{T})\) 的某个有界子集的闭包；
\item
  \(\mathcal{L}^1\big((\mu_\iota)_{\iota \in \mathrm{I}}\big)\) 的每个元素都属于 \(\mathcal{L}_0\) 的某个有界子集的闭包。
\end{enumerate}

为证 a)，可以只考虑形如 \(f = g\varphi_{\mathrm{K}}\) 的元素（ \(g \in \mathcal{C}(\mathrm{T})\) ， K 为 T 中的紧集）。立即看出（借助 Urysohn 定理）， \(f\) 属于由形如 \(gh\) 的函数组成的集 B 的闭包，其中 \(h\) 跑遍 T 到 {[}0, 1{]} 中的连续映射、在 K 上等于 1 且在 K 的一个固定紧邻域 H 的余集上等于 0 者。此外，集 B 是有界的，因为 \(\mu_{\iota}(|gh|) \leqslant \mu_{\iota}(|f\varphi_{\mathrm{H}}|)\) （对每个具有前述性质的函数 h）。

现在来证 b)；可以只考虑形如 \(f \geqslant 0\) 的元素（ \(\mathcal{L}^1((\mu_\iota)_{\iota \in \mathrm{I}})\) 中的）。对每个 \(\iota \in \mathrm{I}\) 与每个 \(\varepsilon > 0\) ，存在一个紧子集 \(\mathrm{K}(\iota, \varepsilon)\) （ \(\mathrm{T}\) 的），使得 \(f\) 在 \(\mathrm{K}(\iota, \varepsilon)\) 上的限制连续且 \(|\mu_\iota(|f - f\varphi_{\mathrm{K}(\iota, \varepsilon)}|) \leqslant \varepsilon\) 。显然 \(f\) 属于集 \(\mathrm{C}\) ——由诸 \(f\varphi_{\mathrm{K}(\iota, \varepsilon)}\) （其中 \(\iota \in \mathrm{I}, \varepsilon > 0\) ）组成——的闭包。借助 Urysohn 定理，集 \(\mathrm{C}\) 含于 \(\mathcal{L}_0\) ；此外它是有界的，因为 \(\mu_\varkappa(f\varphi_{\mathrm{K}(\iota, \varepsilon)}) \leqslant \mu_\varkappa(f)\) （对所有 \(\iota \in \mathrm{I}, \varkappa \in \mathrm{I}\) 与 \(\varepsilon > 0\) ）。

现在来证命题 6：对每个函数 \(g \in \mathcal{K}(\mathrm{T})\) 与每个连续半范数 \(q\) （ \(\mathrm{F}\) 上的），有 \(q\big(\int g dm\big) \leqslant \int |g| d\big(q(\mathbf{m})\big)\) （命题 5），这蕴含映射 \(g \mapsto \int g dm\) （ \(\mathcal{K}(\mathrm{T})\) 到 \(\mathrm{F}\) 中的）当 \(\mathcal{K}(\mathrm{T})\) 赋以 \(\mathcal{L}^1\big((q(\mathbf{m}))_{q \in \mathbb{Q}}\big)\) 的拓扑诱导的拓扑时（Q 为 \(\mathrm{F}\) 上连续半范数之集）是连续的。因此，由前述引理与 TVS, III, §1, No. 6 的命题 10，这一映射可以连续地延拓，先到 \(v_0\) —— \(\mathcal{L}_0\) 到 \(\mathrm{F}\) 中的连续线性映射——再到 \(v\) —— \(\mathcal{L}^1\big((q(\mathbf{m}))_{q \in \mathbb{Q}}\big)\) 到 \(\mathrm{F}\) 中的连续线性映射。此外，对每个 \(\mathbf{z}' \in \mathrm{F}'\) ，关系式 \(\langle \mathbf{z}', v(f) \rangle = \int f d(\mathbf{z}' \circ \mathbf{m})\) （按 \(v\) 的定义）对每个 \(f \in \mathcal{K}(\mathrm{T})\) 成立；由于 \(|\mathbf{z}' \circ \mathbf{m}| \leqslant q(\mathbf{m})\) （对 \(q(\mathbf{z}) = |\langle \mathbf{z}', \mathbf{z} \rangle|\) 而言），映射 \(f \mapsto \int f d(\mathbf{z}' \circ \mathbf{m})\) 在 \(\mathcal{L}^1\big((q(\mathbf{m}))_{q \in \mathbb{Q}}\big)\) 上连续，因此再由连续性，关系式 \(\langle \mathbf{z}', v(f) = \int f d(\mathbf{z}' \circ \mathbf{m})\) 对每个函数 \(f \in \mathcal{L}^1\big((q(\mathbf{m}))_{q \in \mathbb{Q}}\big)\) 成立。由此得出 \(\nu(f) = \int f dm\) ，证明完成。

\subsection{\texorpdfstring{以 \(\mu\) 为基的向量测度}{以 mu 为基的向量测度}}

定义 4. --- 设 \(\mu\) 是 T 上的正测度。称 T 上取值于 F 的向量测度 m 为以 \(\mu\) 为基的测度，如果存在 T 到 F 中的映射 f ，标量局部 \(\mu\)-可积，使得 \(\mathbf{m}(g) = \int g\mathbf{f}d\mu\) （对每个函数 \(g \in \mathcal{K}(\mathrm{T})\) ）。这时称 f 为 m 关于 \(\mu\) 的密度，并记 \(\mathbf{m} = \mathbf{f} \cdot \mu\) 。

立即看出，若 \(f_{1}\) 与 \(f_{2}\) 是 m 关于 \(\mu\) 的两个密度，则 \(f_{1}-f_{2}\) 标量局部 \(\mu\)-可忽略（Ch. V, §5, No. 3, 命题 3 的推论 2）；回忆一般说来这并不蕴含 \(f_{1}-f_{2}\) 局部几乎处处为零（参见 §1, 习题 12 与第 1 目的注记 2）。

设 \(\mathbf{m}\) 是以 \(\mu\) 为基、以 \(\mathbf{f}\) 为密度的测度。为使数值函数 \(g\) 关于 \(\mathbf{m}\) 本质可积，必须且只需 \(gf\) 标量本质 \(\mu\)-可积（Ch. V, §5, No. 3, 定理 1）。

命题 7. --- 设 f 是关于 T 上正测度 \(\mu\) 标量局部可积的映射，使得对每个函数 \(g \in \mathcal{K}(T)\) 有 \(\int g\mathbf{f} d\mu \in F\) 。则映射 \(g \mapsto \int g\mathbf{f} d\mu\) （ \(\mathcal{K}(T)\) 到 F 中的）是 T 上的向量测度，以 \(\mu\) 为基、以 f 为关于 \(\mu\) 的密度。

因为（第 1 目注记 3），只需证明：令 \(\mathbf{m}(g) = \int g\mathbf{f}d\mu\) ，则 \(\mathbf{z}' \circ \mathbf{m}\) 对每个 \(\mathbf{z}' \in \mathrm{F}'\) 是标量测度。但由于 \(\mathbf{z}'(\mathbf{m}(g)) = \int g\langle \mathbf{z}', \mathbf{f} \rangle d\mu\) ，有 \(\mathbf{z}' \circ \mathbf{m} = \langle \mathbf{z}', \mathbf{f} \rangle \cdot \mu\) ，由此即得我们的断言。

命题 8. --- 设 \(\mu\) 是 T 上的正测度， m 是 T 上取值于 F 的测度，以 \(\mu\) 为基、以 f 为关于 \(\mu\) 的密度。设 q 是 F 上的下半连续半范数。

\begin{enumerate}
\def\labelenumi{\alph{enumi})}
\item
  若对每个紧子集 \(\mathrm{K}\) （ \(\mathrm{T}\) 的），上积分 \(\int_{\mathrm{K}}^{*}(q \circ \mathbf{f}) d\mu\) 有限，则 \(\mathbf{m}\) 是 \(q\)-可优超的。
\item
  若 \(\mathbf{m}\) 是 \(q\)-可优超的，则 \(q(\mathbf{m})\) 以 \(\mu\) 为基；若进而 \(\mathbf{f}\) 是 \(\mu\)-可测的（作为 \(\mathrm{T}\) 到 \(\mathrm{F}_{\sigma}\) 中的映射），则 \(q \circ \mathbf{f}\) 局部 \(\mu\)-可积且 \(q(\mathbf{m}) = (q \circ \mathbf{f}) \cdot \mu\) 。
\end{enumerate}

\begin{enumerate}
\def\labelenumi{\alph{enumi})}
\tightlist
\item
  对每个有限子集 \(J\) （ \(A_q'\) 的），用 \(\lambda_J\) 表示诸测度 \(|\mathbf{z}' \circ \mathbf{m}|\) （其中 \(\mathbf{z}'\) 跑遍 \(J\) ）的上确界；若 \(g_J = \sup_{\mathbf{z}' \in J} |\langle \mathbf{z}', \mathbf{f} \rangle|\) ，则 \(\lambda_{J}=g_{J}\cdot\mu\) （Ch. V, §5, No. 2, 命题 2）。对 T 的每个相对紧开子集 U ，令 \(\lambda_{J,U}\) 为 \(\lambda_{J}\) 在 U 上的限制；首先证明：当 J 跑遍 \(A_{q}^{\prime}\) 的有限子集组成的有向集 F 时，族 \((\lambda_{J,U})\) 在 \(\mathcal{M}(U)\) 中有上界。的确，对每个函数 \(h\geqslant0\) （ \(K(U)\) 中的），
\end{enumerate}

\[
\int h d \lambda_ {\mathrm{J}, \mathrm{U}} = \int h g _ {\mathrm{J}} d \mu \leqslant \int^ {*} (q \circ \mathbf {f}) h d \mu \leqslant \| h \| \int_ {\mathrm{U}} ^ {*} (q \circ \mathbf {f}) d \mu ,
\]

由此即得我们的断言（Ch. II, §2, No. 2）。令 \(\nu_{\mathrm{U}}\) 为这族测度在 \(\mathcal{M}(\mathrm{U})\) 中的上确界。若 \(\mathrm{U}'\) 是 T 的另一个相对紧开子集，满足 \(\mathrm{U} \subset \mathrm{U}'\) ，则 \(\nu_{\mathrm{U}}\) 是 \(\nu_{\mathrm{U}'}\) 在 U 上的限制，这立即由递增有向测度集的上确界的表达式（Ch. II, §2, No. 2）以及 \(\lambda_{\mathrm{J,U}}\) 是 \(\lambda_{\mathrm{J,U}'}\) 到 U 上的限制这一事实得出。于是存在唯一的正测度 \(\nu\) ，它在每个 U 上的限制是 \(\nu_{\mathrm{U}}\) （Ch. III, §2, No. 1, 命题 1），且显然 \(\nu = q(\mathbf{m})\) 。

\begin{enumerate}
\def\labelenumi{\alph{enumi})}
\setcounter{enumi}{1}
\tightlist
\item
  由于诸测度 \(\lambda_{\mathrm{J}}\) 以 \(\mu\) 为基，其上确界 \(q(\mathbf{m})\) 亦然（Ch. V, §5, No. 5, 定理 2）。若 \(\mathbf{f}\) 是 \(\mu\)-可测的（对拓扑 \(\sigma(\mathrm{F},\mathrm{F}')\) ， \(\mathrm{F}\) 上的），则由定义立即得出，映射 \(g:t\mapsto(g_{\mathrm{J}}(t))_{\mathrm{J}\in\mathfrak{F}}\) （ \(\mathrm{T}\) 到乘积空间 \(\mathbf{R}^{\mathfrak{F}}\) 中的）是 \(\mu\)-可测的。 \(q\circ\mathbf{f}=\sup_{\mathrm{J}\in\mathfrak{F}}g_{\mathrm{J}}\) 到 \(\mathrm{T}\) 的每个使 \(g\) 连续的紧子集上的限制是下半连续的；从而 \(q \circ f\) 是 \(\mu\)-可测的（Ch. IV, §5, No. 5, 命题 8 的推论及 No. 10, 命题 15）。设 \(K\) 是 T 的紧子集；它容许一个分割，由一个 \(\mu\)-可忽略集与 \((\mathrm{K}_n)\) （ \(g\) 在其上连续的紧集序列）组成。这时 \(\int_{\mathrm{K}_n}^*\left(q\circ \mathbf{f}\right)d\mu = \sup_{\mathrm{J}}\int_{\mathrm{K}_n}g_{\mathrm{J}}d\mu \leqslant \int_{\mathrm{K}_n}dq(\mathbf{m})\) （对所有 \(n\) ）（Ch. IV, §1, No. 1, 定理 1 与 Ch. V, §7, No. 1, 命题 1），由此 \(\int_{\mathrm{K}}^{*}(q\circ \mathbf{f})d\mu = \sum_{n}\int_{\mathrm{K}_n}^*\left(q\circ \mathbf{f}\right)d\mu \leqslant \int_{\mathrm{K}}dq(\mathbf{m})\) 。但这证明了 \(q\circ \mathbf{f}\) 局部 \(\mu\)-可积，且 \(\lambda_{\mathrm{J}}\leqslant (q\circ \mathbf{f})\cdot \mu \leqslant q(\mathbf{m})\) （对所有 \(\mathbf{J}\in \mathfrak{F}\) ）；于是按定义， \(q(\mathbf{m}) = (q\circ \mathbf{f})\cdot \mu\) 。
\end{enumerate}

注记. --- 假设 \(A_{q}^{\prime}\) 中存在对 \(\sigma(\mathbf{F}^{\prime},\mathbf{F})\) 稠密的可数子集 D ；那么函数 \(q\circ f\) 总是 \(\mu\)-可测的，因为 \(q(\mathbf{f}(t)) = \sup_{\mathbf{z}^{\prime}\in\mathrm{D}} |\langle\mathbf{z}^{\prime},\mathbf{f}(t)\rangle|\) （Ch. IV, §5, No. 4, 定理 2 的推论 1）。于是，对

T 的每个紧子集 K ，有 \(\int_{K}^{*}(q\circ\mathbf{f})d\mu=\sup_{J}\int_{K}g_{J}d\mu\) ，其中 J 跑遍 D 的有限子集组成的可数有向集（Ch. IV, §1, No. 1, 定理 3 的推论）；于是可以看出，在这种情形，条件 \(\int_{K}^{*}(q\circ\mathbf{f})d\mu<+\infty\) （对 T 的每个紧子集 K ）是 m 为 q-可优超的必要充分条件。

命题 9. --- 设 F 是有限维 Banach 空间。T 上取值于 F 的每个测度 m 都是以 \textbar m\textbar 为基的测度。若 \(m = f \cdot |m|\) ，则 \(|f(t)| = 1\) 局部几乎处处成立（关于 \textbar m\textbar ）。为使 m 以 \(\mu\) 为基，必须且只需 \textbar m\textbar{} 以 \(\mu\) 为基，且若 \(m = g \cdot \mu\) 则 \(|m| = |g| \cdot \mu\) 。

设 \((\mathbf{e}_i)_{1 \leqslant i \leqslant n}\) 与 \((\mathbf{e}_i')_{1 \leqslant i \leqslant n}\) 是 F 与 \(\mathbf{F}'\) 的对偶基（A, II, §2, No. 6），满足 \(|\mathbf{e}_i'| = 1\) （对所有 \(i\) ）。这时 \(|\mathbf{e}_i' \circ \mathbf{m}| \leqslant |\mathbf{m}|\) （对每个指标 \(i\) ），因此（Ch. V, §5, No. 5, 定理 2） \(\mathbf{e}_i' \circ \mathbf{m} = h_i \cdot |\mathbf{m}|\) ，其中 \(h_i\) 有界且 \(|\mathbf{m}|\)-可测。令 \(\mathbf{h} = \sum_{i=1}^{n} h_i \cdot \mathbf{e}_i\) ，我们便有 \(\mathbf{m} = \mathbf{h} \cdot |\mathbf{m}|\) 。若 \(\mathbf{m} = \mathbf{f} \cdot |\mathbf{m}|\) ，命题 8 表明 \(|\mathbf{m}| = |\mathbf{f}| \cdot |\mathbf{m}|\) ，从而 \(|\mathbf{f}(t)| = 1\) 关于 \(|\mathbf{m}|\) 局部几乎处处成立（Ch. V, §5, No. 3, 命题 3 的推论 2）。最后的断言立即由命题 8 得出。

注记. --- 若 \(z = \sum_{i=1}^{n} z_i e_i\) ，则 \(\psi(z_1, \ldots, z_n) = |z|\) 是 \(R to the n\) 上的正齐次连续函数。令 \(\mu_i = e_i' \circ m = h_i \cdot |m|\) ，按定义（Ch. V, §5, No. 9）有 \(\psi(\mu_1, \ldots, \mu_n) = \psi(h_1, \ldots, h_n) \cdot |m| = |h| \cdot |m| = |m|\) 。

\subsection{Dunford--Pettis 定理}

设 \(\mu\) 是 T 上的正测度。称 T 上取值于 F 的向量测度 \(\mathbf{m}\) 为标量地以 \(\mu\) 为基（或标量地以 \(\mu\) 为基），如果对每个 \(\mathbf{z}' \in \mathrm{F}'\) ，标量测度 \(\mathbf{z}' \circ \mathbf{m}\) 以 \(\mu\) 为基。若取值于 F 的向量测度 m 以 \(\mu\) 为基，则它标量地以 \(\mu\) 为基：因为若 \(m = f \cdot \mu\) ，则 \(z' \circ m = \langle z', f \rangle \cdot \mu\) （对每个 \(z' \in F'\) ）。但存在标量地以 \(\mu\) 为基而不以 \(\mu\) 为基的向量测度（习题 17），另一方面，也存在对任何正测度 \(\nu\) 都不标量地以 \(\nu\) 为基的向量测度；但注意，取值于赋范空间的每个可优超测度 m 依 Lebesgue--Nikodym 定理都标量地以 \(|m|\) 为基。

例. --- 取 m 为 \(\mathcal{K}(\mathrm{T})\) 的恒等映射。说 m 标量地以 \(\mu\) 为基，就是说 T 上每个实测度都以 \(\mu\) 为基。特别地，点测度 \(\varepsilon_{t}\) ( \(t \in T\) ) 必须以 \(\mu\) 为基，这要求 \(\mu(\{t\}) > 0\) （对每个 \(t \in T\) ），并特别蕴含 T 的每个紧子集都可数。

在本目中，我们将证明一个结果，它推广了 Lebesgue--Nikodym 定理的推论之一，即 \(\mathrm{L}^{1}(\mu)\) 的对偶是 \(\mathrm{L}^{\infty}(\mu)\) （Ch. V, §5, No. 8, 定理 4），并且给出标量地以 \(\mu\) 为基的向量测度以 \(\mu\) 为基的一个充分条件。

设 \(\pi\) 是 \(\mathcal{L}^{\infty}(\mu)\) 到 \(\mathrm{L}^{\infty}(\mu)\) 上的典范映射。称线性子空间 \(\mathbf{G}\) （ \(\mathrm{L}^{\infty}\) 的）具有提升性质，如果存在线性映射 \(\rho\) （ \(\mathbf{G}\) 到 \(\mathcal{L}^{\infty}(\mu)\) 中的；称为 \(\mathbf{G}\) 的提升），使得 \(\pi \circ \rho\) 是 \(\mathbf{G}\) 上的恒等映射，且 \(|\rho(f)(t)| \leqslant \mathrm{N}_{\infty}(f)\) （对所有 \(t \in \mathrm{T}\) 与 \(f \in \mathrm{G}\) ）。

可以证明，若 \(\mu\) 是 \(\mathbf{R}^n\) 上的 Lebesgue 测度，则整个空间 \(\mathrm{L}^{\infty}(\mathbf{R}^{n},\mu)\) 具有提升性质（习题 18）。

引理 2. --- Banach 空间 \(\mathrm{L}^{\infty}(\mathrm{T},\mu)\) 的每个可分子空间 G 都具有提升性质。

按假设，存在 G 的可数稠密子集 H ，它是关于有理数域 Q 的线性子空间；设 \((h_{n})\) 是 H 在 Q 上的一个（可数）基。对每个整数 n ，令 \(h_{n}^{\prime}\) 为 \(L^{\infty}\) 中满足 \(\pi(h_{n}^{\prime}) = h_{n}\) 的元素，令 \(\rho^{\prime}\) 为 H 到 \(L^{\infty}\) 中的由 \(\rho^{\prime}(h_{n}) = h_{n}^{\prime}\) 定义的 Q-线性映射；显然 \(\pi \circ \rho^{\prime}\) 是 H 上的恒等映射。此外，对每个 \(h \in H\) 有 \(|\rho^{\prime}(h)(t)| \leqslant N_{\infty}(h)\) ，但在局部可忽略集 A(h) 的点 t 处除外。令 A 为诸 A(h) （ \(h \in H\) ）的并，它也是局部可忽略的。对每个 \(h \in H\) ，用 \(\rho(h)\) 表示满足如下条件的函数 \(h^{\prime\prime} \in L^{\infty}\) ： \(h^{\prime\prime}(t) = \rho^{\prime}(h)(t)\) （若 \(t \notin A\) ），而 \(h^{\prime\prime}(t) = 0\) （若 \(t \in A\) ）。显然 \(\rho\) 是 H 到 \(L^{\infty}\) 中有界函数所成的子空间 B 中的 Q-线性映射，使得 \(\pi \circ \rho\) 是 H 上的恒等映射，且 \(|\rho(h)(t)| \leqslant N_{\infty}(h)\) （对所有 \(h \in H\) 与 \(t \in T\) ）。由于 B 对范数 \(\|f\| = \sup_{t \in T} |f(t)|\) （Ch. IV, §5, No. 4, 定理 2）是 Banach 空间，

\(\rho\) 可以延拓为一个连续的 R-线性映射（仍记作 \(\rho\) ，由 G 到 B 中），它显然是 G 的一个提升。

定义 5. --- 设 F 是 Hausdorff 局部凸空间， \(F_{s}^{\prime}\) 是它的对偶，赋以拓扑 \(\sigma(F^{\prime},F)\) 。我们用 \(L_{F_{s}^{\prime}}^{\infty}\) 表示 T 到 \(F_{s}^{\prime}\) 中的映射 f 所成的向量空间，其中 f 标量 \(\mu\)-可测，且标量地局部几乎处处（对 \(\mu\) ）等于一个 T 到 \(F^{\prime}\) 的等度连续子集中的映射。我们用 \(L_{F_{s}^{\prime}}^{\infty}\) 表示 \(L_{F_{s}^{\prime}}^{\infty}\) 关于标量局部 \(\mu\)-可忽略的 T 到 \(F_{s}^{\prime}\) 中的映射所成空间的商空间。

当 F 满足 §1 第 5 目命题 13 的假设时， \(\mathcal{L}_{\mathrm{F}_s^\prime}^\infty\) 中的函数是 \(\mu\)-可测的（对弱拓扑 \(\sigma(\mathrm{F}',\mathrm{F})\) 而言），但对 \(\mathrm{F}'\) 上的强拓扑不必可测，即使 F 是 Banach 空间（§1, 习题 17）。在同样条件下，标量局部 \(\mu\)-可忽略的 T 到 \(\mathrm{F}_s^\prime\) 中的映射与局部 \(\mu\)-可忽略的 T 到 \(\mathrm{F}_s^\prime\) 中的映射相同（§1, 第 1 目的注记 2）。

当 F 是可分赋范空间时， \(L_{F_{s}^{\prime}}^{\infty}\) 的元素是 T 到 \(F_{s}^{\prime}\) 中的映射 f ，其中 f 标量 \(\mu\)-可测且 \(|f|\) 依测度有界；这时可以在空间 \(L_{F_{s}^{\prime}}^{\infty}\) 上定义赋范空间结构，即赋以范数 \(N_{\infty}\) （Ch. IV, §6, No. 3）。

引理 3. --- 设 F 是 Hausdorff 局部凸空间， f 是 \(\mathcal{L}_{\mathrm{F}_s^\prime}^\infty\) 的元素。对每个 \(\mathbf{z} \in \mathrm{F}\) ，有 \(\langle \mathbf{z}, \mathbf{f} \rangle \in \mathcal{L}^\infty\) ，且线性映射 \(z \mapsto \pi(\langle \mathbf{z}, \mathbf{f} \rangle)\) （F 到 \(\mathrm{L}^\infty\) 中的）连续；若进而 F 是赋范空间，则 \(\mathrm{N}_\infty(\langle \mathbf{z}, \mathbf{f} \rangle) \leqslant |\mathbf{z}| \cdot \sup_{t \in \mathrm{T}} |\mathbf{f}(t)|\) 。

按定义显然 \(\langle z,f\rangle\) 是 \(\mu\)-可测的且依测度有界；若有必要，把 f 换成属于 \(L_{F_{s}}^{\infty}\) 中同一类的函数，则可以进而假设 \(\mathbf{f}(T)\subset V^{\circ}\) ，其中 V 是 F 中 0 的平衡凸邻域（这在至多一个局部可忽略集上不改变 \(\langle z,f\rangle\) ，该集依赖于 z）。这时关系式 \(z\in V\) 蕴含 \(|\langle z,f(t)\rangle|\leqslant1\) （对所有 \(t\in T\) ），这就证明了 \(\mathbf{z}\mapsto\pi(\langle\mathbf{z},\mathbf{f}\rangle)\) 的连续性。最后的断言是显然的。

引理 4. --- 设 F 是 Hausdorff 局部凸空间， f 是 \(\mathcal{L}_{\mathrm{F}_s}^\infty\) 的元素。对每个数值函数 \(g \in \overline{\mathcal{L}}^1\) ，函数 gf 标量本质 \(\mu\)-可积，且 \(\int g\mathbf{f} d\mu \in \mathrm{F}'\) 。

对每个 \(z \in F\) ， \(\langle z, f \rangle\) 属于 \(L^{\infty}\) ，由此得第一个断言。可以不改变 \(\int g f d\mu\) 而进而假设 \(\mathbf{f}(T) \subset V^{\circ}\) ，其中 V 是 F 中 0 的平衡凸邻域。这时关系式 \(z \in V\) 蕴含 \(|\langle z, f(t) \rangle| \leqslant 1\) （对所有 \(t \in T\) ），从而 \(|\langle z, \int g f d\mu \rangle| = |\int \langle z, f \rangle g d\mu| \leqslant \overline{N}_{1}(g)\) ，这证明了 \(\int g f d\mu \in F'\) 。

定理 1. --- 设 F 是包含可数稠密集的 Hausdorff 局部凸空间。对每个函数 \(\mathbf{f} \in \mathcal{L}_{\mathrm{F}_s}^\infty\) 与每个 \(\mathbf{z} \in \mathbf{F}\) ，令 \(v_{\mathbf{f}}(\mathbf{z}) = \pi (\langle \mathbf{z}, \mathbf{f} \rangle) \in \mathrm{L}^{\infty}\) ；映射 \(\mathbf{f} \mapsto v_{\mathbf{f}}\) 经取商定义了 \(L_{F^{\prime}s}^{\infty}\) 到空间 \(\mathcal{L}(F;L^{\infty})\) （F 到 \(L^{\infty}\) 中的连续线性映射所成的空间）上的线性双射。若 F 是赋范空间，则此双射是等距的。

鉴于引理 3，若能证明：对每个连续映射 \(u\) （ \(F\) 到 \(L^{\infty}\) 中的），存在函数 \(\mathbf{f} \in \mathcal{L}_{F_s'}^{\infty}\) 使得 \(\pi(\langle \mathbf{z}, \mathbf{f} \rangle) = u(\mathbf{z})\) （对所有 \(\mathbf{z} \in F\) ），且 \(\mathbf{f}\) 在 \(L_{F_s'}^{\infty}\) 中的类由该条件唯一确定，则第一个断言得证。第二点是直接的，因为若 \(\pi(\langle \mathbf{z}, \mathbf{f} \rangle) = \pi(\langle \mathbf{z}, \mathbf{f}_1 \rangle)\) （对所有 \(\mathbf{z} \in F\) ），则 \(\mathbf{f}_1 - \mathbf{f}\) 标量局部可忽略。另一方面，存在提升 \(\rho\) （ \(u(F)\) 到 \(\mathcal{L}^{\infty}\) 中的）（引理 2）。对每个 \(t \in T\) ，映射 \(\mathbf{z} \mapsto \rho(u(\mathbf{z}))(t)\) 是 \(F\) 上的连续线性形式，即一个元素 \(f(t)\) （ \(F'\) 的）。函数 \(\mathbf{f}\) 标量 \(\mu\)-可测，因为 \(\langle \mathbf{z}, \mathbf{f} \rangle = \rho(u(\mathbf{z})) \in \mathcal{L}^{\infty}\) （对每个 \(\mathbf{z} \in F\) ）；有 \(\pi(\langle \mathbf{z}, \mathbf{f} \rangle) = u(\mathbf{z})\) ；最后，对每个 \(t \in T\) 与每个 \(\mathbf{z}\) ——属于逆像 \(V\) （ \(u\) 下 \(L^{\infty}\) 的单位球的）者——，

\[
\left| \langle \mathbf {z}, \mathbf {f} (t) \rangle \right| = \left| \rho (u (\mathbf {z})) (t) \right| \leqslant N _ {\infty} (u (\mathbf {z})) \leqslant 1,
\]

这表明 \(\mathbf{f}(t) \in \mathrm{V}^{\circ}\) （对所有 \(t \in \mathrm{T}\) ）。

若进而 F 是赋范空间，前面的论证表明

\[
\sup _ {t \in \mathrm{T}} | \mathbf {f} (t) | \leqslant \| u \|.
\]

但另一方面（引理 3）， \(\mathrm{N}_{\infty}(u(\mathbf{z})) \leqslant |\mathbf{z}| \cdot \sup_{t \in \mathrm{T}} |\mathbf{f}(t)|\) ，且当 f 在一个局部可忽略集上被任意修改时，该不等式继续成立。由此得出 \(\|u\| = \mathrm{N}_{\infty}(|\mathbf{f}|)\) 。

推论 1. --- 设 F 是包含可数稠密集的 Hausdorff 局部凸空间。对每个函数 \(\mathbf{f} \in \mathcal{L}_{\mathrm{F}_s}^\infty\) 、每个 \(\mathbf{z} \in \mathrm{F}\) 与每个函数 \(g \in \mathcal{L}^1\) ，令 \(\Phi_{\mathbf{f}}(\mathbf{z}, \widetilde{g}) = \int \langle \mathbf{z}, \mathbf{f}(t) \rangle g(t) d\mu(t)\) 。映射 \(\mathbf{f} \mapsto \Phi_{\mathbf{f}}\) 经取商定义了 \(\mathrm{L}_{\mathrm{F}_s}^\infty\) 到空间 \(\mathcal{B}(\mathrm{F}, \mathrm{L}^1)\) （ \(\mathrm{F} \times \mathrm{L}^1\) 上连续双线性形式所成的空间）上的线性双射。若 \(\mathrm{F}\) 是赋范空间，则此双射是等距的。

可以假设 \(\mathbf{f}(\mathrm{T})\) 是 \(F'\) 的等度连续子集。这时显然 \(\Phi_{f}\) 分别连续，且用定理 1 与附录的记号，有（考虑到 \(L^{\infty}\) 是 \(L^{1}\) 的对偶这一事实（Ch. V, §5, No. 8, 定理 4）） \(^{l}\Phi_{f}=v_{f}\) 。于是推论由上面的定理 1 以及附录第 1 目命题 1 及其推论得出。

推论 2（Dunford--Pettis 定理）. --- 设 F 是包含可数稠密集的 Hausdorff 局部凸空间。对每个函数 \(f \in L_{F_{s}^{\infty}}\) 与每个函数 \(g \in L^{1}\) ，令 \(w_{\mathbf{f}}(\widetilde{g}) = \int g\mathbf{f} \, d\mu \in \mathbf{F}'\) （引理 4）。映射 \(f \mapsto w_{f}\) 经取商定义了 \(L_{F_{s}^{\prime}}^{\infty}\) 到空间 \(\mathcal{R}(L^{1}, F^{\prime})\) 上的线性双射，该空间由 \(L^{1}\) 到 \(F^{\prime}\) 中的线性映射 u 组成，使得 u 下 \(L^{1}\) 的单位球的像是 \(F^{\prime}\) 的等度连续子集。若 F 是赋范空间（这时 \(\mathcal{R}(L^{1}, F^{\prime})\) 是 \(L^{1}\) 到 F 的强对偶中的连续线性映射所成的空间），则双射 \(f \mapsto w_{f}\) 是等距的。

考虑到 \(L^{\infty}\) 是 \(L^{1}\) 的对偶这一事实，这由前述推论以及附录第 1 目命题 1 及其推论得出。

注记. --- 显然，映射 \(u \in \mathcal{R}(L^{1}, F')\) 对 \(F'\) 上每个 S-拓扑都连续（S 是 F 的由有界子集组成的覆盖）。反之，若进而假设 F 是桶式的，则 \(L^{1}\) 到赋以某 S-拓扑的 \(F'\) 中的每个连续线性映射都把 \(L^{1}\) 的单位球变成 \(F'\) 的有界子集，从而是等度连续的（TVS, III, §4, No. 2, 定理 1）。

推论 3. --- 设 F 是包含可数稠密集的 Hausdorff 局部凸空间， m 是 T 上取值于 F 的弱对偶 F' 的向量测度。若函数集 B （ \(\mathcal{K}(T)\) 中的函数 g ，满足 \(\mu(|g|) \leqslant 1\) 者）在 m 下的像含于 F' 的一个闭凸等度连续子集 H' ，则 m 以 \(\mu\) 为基，且存在 m 关于 \(\mu\) 的密度 f ，使得对所有 t ∈ T 有 f(t) ∈ H'。

假设蕴含：当 \(\mathcal{K}(T)\) 赋以 \(\mathcal{L}^{1}(\mu)\) 的拓扑诱导的拓扑（由半范数 \(N_{1}\) 定义）时，m 连续；因此它可延拓为 \(\mathcal{L}^{1}(\mu)\) 到 F 的弱对偶的完备化 G 中的连续线性映射 u；但由 \(H'\) 是 G 的紧子集且集合 \(\overline{B}\) ——由 \(f \in L^{1}\) （满足 \(\mathrm{N}_{1}(f) \leqslant 1\) 者）组成——在 u 下的像含于 \(H'\) 在 G 中的闭包，有 \(u(\overline{\mathrm{B}}) \subset \mathrm{H}'\) ，因此 u 把 \(L^{1}\) 映入 \(F'\) 。由于关系式 \(\mathrm{N}_{1}(f) \leqslant \varepsilon\) 蕴含 \(u(f) \in \varepsilon\mathrm{H}'\) ，故有 \(u(g) = 0\) （若 g \(\mu\)-可忽略），于是推论 2 可应用于由 u 经取商得到的 \(L^{1}\) 到 \(F'\) 中的映射；由此即得推论。

推论 4. --- 设 F 是可分赋范空间， m 是 T 上取值于 F 的强对偶 F' 的测度，对 F' 的范数可优超。则 m 是以 \textbar m\textbar 为基的测度，且若 \(m = f \cdot |m|\) ，则 \(|f(t)| = 1\) 关于 \textbar m\textbar 局部几乎处处成立。

按假设，对每个 \(\mathbf{z} \in \mathbf{F}\) （满足 \(|\mathbf{z}| \leqslant 1\) 者），有 \(|\langle \mathbf{z}, \mathbf{m}(g) \rangle| \leqslant |\mathbf{m}|(|g|)\) （对每个函数 \(g \in \mathcal{K}(\mathrm{T})\) ），从而 \(|\mathbf{m}(g)| \leqslant |\mathbf{m}|(|g|)\) （TVS, IV, §2, No. 4, 公式 (1)）。由于 \(\mathbf{F}'\) 中每个球都等度连续，推论 3 适用并表明 \(\mathbf{m}\) 以 \(|\mathbf{m}|\) 为基；此外，若 \(\mathbf{m} = \mathbf{f} \cdot |\mathbf{m}|\) ，则 \(\mathbf{f}\) 是 \(|\mathbf{m}|\)-可测的（对 \(\sigma(\mathrm{F}', \mathrm{F})\) 而言）（§1, 第 5 目，命题 13），且 \(|\mathbf{m}| = |\mathbf{f}| \cdot |\mathbf{m}|\) （第 4 目，命题 8），这就证明了推论（Ch. V, §5, No. 3, 命题 3 的推论 2）。

若把这推论应用于 \(\mathbf{F}\) 为有限维的情形，就重新得到命题 9 的第一部分作为特例。

\subsection{\texorpdfstring{空间 \(L_{F}^{1}\) 的对偶（F 为可分 Banach 空间）}{空间 L sub F superscript 1 的对偶（F 为可分 Banach 空间）}}

命题 10. --- 设 F 是可分 Banach 空间。对每个函数 \(\mathbf{f} \in \overline{\mathcal{L}}_{\mathrm{F}}^{1}\) 与每个函数 \(\mathbf{g} \in \mathcal{L}_{\mathrm{F}'_s}^{\infty}\) ，数值函数 \(\langle \mathbf{f}, \mathbf{g} \rangle : t \mapsto \langle \mathbf{f}(t), \mathbf{g}(t) \rangle\) 本质 \(\mu\)-可积，且

\[
\left| \int \langle \mathbf {f}, \mathbf {g} \rangle d \mu \right| \leqslant \overline {{{{N}}}} _ {1} (\mathbf {f}) N _ {\infty} (\mathbf {g}).\tag{3}
\]

对每个类 \(\dot{\mathbf{g}}\in\mathrm{L}_{\mathrm{F}^{\prime}s}^{\infty}\) ，令 \(\theta(\dot{\mathbf{g}})\) 为 \(\mathrm{L}_{\mathrm{F}}^{1}\) 上的连续线性形式，它是由线性形式 \(\mathbf{f}\mapsto\int\langle\mathbf{f},\mathbf{g}\rangle d\mu\) （ \(\overline{\mathcal{L}}_{\mathrm{F}}^{1}\) 上的）经取商得到的；这时 \(\theta\) 是 \(\mathrm{L}_{\mathrm{F}^{\prime}s}^{\infty}\) 到强对偶 \((\mathrm{L}_{\mathrm{F}}^{1})^{\prime}\) （Banach 空间 \(\mathrm{L}_{\mathrm{F}}^{1}\) 的）上的线性等距。

对 T 的每个紧子集 K 与每个 \(\varepsilon > 0\) ，存在 K 的紧子集 \(K'\) ，使得 \(\mu(K - K') \leqslant \varepsilon\) ，且 f（相应地， g）到 \(K'\) 的限制是 \(K'\) 到 F（相应地，到 \(F'_s\) ）中的连续映射；因此 \(\mathbf{g}(K')\) 是弱紧的，从而在 \(F'\) 上等度连续（TVS, III, §4, No. 2, 定理 1 或 IV, §1, 习题 10）。而今， \(F \times F'\) 上的典范双线性形式到 F 与 \(F'\) 的等度连续子集的乘积上的限制，对 F 的拓扑与 \(\sigma(F', F)\) 的乘积拓扑是连续的（GT, X, §2, No. 1, 命题 1 的推论 4）；由此， \(\langle f, g \rangle\) 到 \(K'\) 的限制连续，从而 \(\langle f, g \rangle\) 是 \(\mu\)-可测的。此外，

\[
| \langle \mathbf {f} (t), \mathbf {g} (t) \rangle | \leqslant | \mathbf {f} (t) | \cdot | \mathbf {g} (t) | \leqslant | \mathbf {f} (t) | N _ {\infty} (\mathbf {g})
\]

局部几乎处处成立，从而 \(\langle f,g\rangle\) 本质 \(\mu\)-可积且不等式 (3) 成立（Ch. IV, §5, No. 6, 定理 5 及 Ch. V, §1, No. 3, 引理）。

剩下要证 \(\theta\) 是满射等距。设 u 是 \(L_{F}^{1}\) 上的连续线性形式。映射 \((h,\mathbf{z})\mapsto u(\widetilde{h}\mathbf{z})\) 是 \(L^{1}\times F\) 上的连续双线性形式，因为

\[
\left| u \left(\widetilde {h} \mathbf {z}\right) \right| \leqslant \| u \| \cdot \mathrm{N} _ {1} (h \mathbf {z}) \leqslant \| u \| \cdot | \mathbf {z} | \cdot \mathrm{N} _ {1} (h).
\]

由第 5 目定理 1 的推论 1，存在 T 到 F' 中的映射 g ，属于 \(L_{F_{s}^{\infty}}\) ，使得 \(|\mathbf{g}(t)| \leqslant \|u\|\) （对所有 \(t \in T\) ），且 \(u(\widetilde{h}\mathbf{z}) = \int \langle h\mathbf{z}, \mathbf{g} \rangle d\mu\) （对每个函数 \(h \in L^{1}\) ——其类 \(\widetilde{h}\) 在 \(L^{1}\) 中——与每个 \(z \in F\) ）。换言之，线性形式 u 与 \(\theta(\dot{\mathbf{g}})\) 在 \(L_{F}^{1}\) 的由形如 \(\widetilde{h}\mathbf{z}\) 的元素（ \(\widetilde{h} \in L^{1}\) , \(z \in F\) ）生成的子空间上一致。由于该子空间在 \(L_{F}^{1}\) 中稠密（Ch. IV, §3, No. 5, 命题 10），得出 \(u = \theta(\dot{\mathbf{g}})\) ，这已经证明 \(\theta\) 是满射的。此外，由 (3)， \(\|\theta(\dot{\mathbf{g}})\| \leqslant N_{\infty}(\mathbf{g}) \leqslant \|u\| = \|\theta(\dot{\mathbf{g}})\|\) ，从而 \(\|\theta(\dot{\mathbf{g}})\| = N_{\infty}(\mathbf{g})\) ，证明完毕。
